\documentclass[twoside,11pt,openright]{report}

% \usepackage[latin1]{inputenc}
\usepackage[english]{babel}
\usepackage{a4}
\usepackage{latexsym}
\usepackage{amssymb}
\usepackage{amsmath}
\usepackage{epsfig}
\usepackage[T1]{fontenc}
% \usepackage{mathptmx} % Use Times New Romans as font?
\usepackage{color}
\usepackage{epstopdf}
\usepackage{microtype}
\usepackage[hidelinks]{hyperref}
\usepackage[useregional]{datetime2}
\DTMlangsetup[en-US]{showdayofmonth=false}
\usepackage{lipsum}

\usepackage{marvosym}

\renewcommand*\sfdefault{lmss}
\renewcommand*\ttdefault{txtt}

% Custom packages added:

\usepackage{fullpage}
\setlength\marginparwidth{0.7in}

\usepackage{todonotes} % \todo command

\usepackage{proof} % \infer command

\usepackage{tikz-cd}
\usepackage{xfrac} % \sfrac - sideways fraction

% allow breaks in align and equation (a new page for example) 
\allowdisplaybreaks

\usepackage{amsthm}

\theoremstyle{plain} % default
\newtheorem{thm}{Theorem}[section]
\newtheorem{lem}[thm]{Lemma}
\newtheorem{prop}[thm]{Proposition}
\newtheorem*{cor}{Corollary}

\theoremstyle{definition}
\newtheorem{defn}[thm]{Definition}% [section] 
\newtheorem{conj}{Conjecture}[section]
\newtheorem{exmp}{Example} % [section]

\theoremstyle{remark}
\newtheorem*{rem}{Remark}
\newtheorem*{note}{Note}
\newtheorem{case}{Case}

% Custom commands

\newcommand*{\thmref}[1]{\hyperref[thm:#1]{Theorem~\ref*{thm:#1}}} % Theorem~\ref
\newcommand*{\lemref}[1]{\hyperref[lem:#1]{Lemma~\ref*{lem:#1}}}
\newcommand*{\figref}[1]{\hyperref[fig:#1]{Figure~\ref*{fig:#1}}}

\newcommand*{\set}[1]{\left\{#1\right\}}
\newcommand*{\term}[1]{\textcolor{green!30!black}{#1}} % \textcolor{green!70!black} \textcolor{blue!80!black}
\newcommand*{\pathterm}[1]{\textcolor{green!40!black}{#1}}
\newcommand*{\type}[1]{\textcolor{magenta!90!black}{#1}}
% \newcommand*{\dependenttype}[1]{\textcolor{magenta!70!black}{#1}}
\newcommand*{\container}[1]{\textcolor{orange}{#1}}
\newcommand*{\containerpair}[2]{\ensuremath{\colorlet{savedleftcolor}{.}\color{orange}\left(\color{savedleftcolor}#1\,\textcolor{orange}{\mathbf{,}}\,#2\color{orange}\right)\color{savedleftcolor}}}
\newcommand*{\containerpairsimple}[2]{\containerpair{#1}{\lambda\,\_,\,#2}}
\newcommand*{\universe}[1]{\textcolor{orange!80!black}{#1}}
\newcommand*{\unit}{\type{\textbf{1}}}
\newcommand*{\empt}{\type{\textbf{0}}}
\newcommand*{\coalg}[2]{#1\texttt{-}#2}

\newcommand*{\relation}[1]{\textcolor{gray!75!black}{\ensuremath{\mathtt{#1}}}}
\newcommand*{\constant}[1]{\textcolor{orange!60!black}{\ensuremath{\mathtt{#1}}}}

\newcommand*{\function}[1]{\textcolor{blue!60!black}{\ensuremath{\mathtt{#1}}}}
\newcommand*{\constructor}[1]{\textcolor{purple!60!black}{\ensuremath{\mathtt{#1}}}}
\newcommand*{\destructor}[1]{\textcolor{yellow!60!black}{\ensuremath{\mathtt{#1}}}}
\newcommand*{\typeformer}[1]{\ensuremath{\mathtt{#1}}}
\newcommand*{\functor}[1]{\ensuremath{\mathbf{\mathtt{#1}}}}

\newcommand*{\unitelem}{\constant{\star}} % \constant{tt}

\newcommand*{\ssfrac}[2]{ {\LARGE \text{\sfrac{\ensuremath{#1}}{\ensuremath{#2}}}}}

\newcommand*{\natcases}[2]{\ensuremath{\mathbb{[\hspace{-.5mm}\rangle}~\ensuremath{#1}~,~\ensuremath{#2}~\mathbb{\langle\hspace{-.5mm}]}}}

% Concatenation of paths taken from HoTT book
\newcommand{\ct}{%
  \mathchoice{\mathbin{\raisebox{0.5ex}{$\displaystyle\centerdot$}}}%
             {\mathbin{\raisebox{0.5ex}{$\centerdot$}}}%
             {\mathbin{\raisebox{0.25ex}{$\scriptstyle\,\centerdot\,$}}}%
             {\mathbin{\raisebox{0.1ex}{$\scriptscriptstyle\,\centerdot\,$}}}
}

\newcommand*{\sym}[1]{\ensuremath{#1^{-1}}}

% \mathord

\begin{document}

%%%%%%%%%%%%%%%%%%%%%%%%%%%%%%%%%%%%%%%%%%%%%%%%%%%%%%%%%%%%%%%%%%%%%%% 

\pagestyle{empty} 
\pagenumbering{roman} 
\vspace*{\fill}\noindent{\rule{\linewidth}{1mm}\\[4ex]
  {\Huge\sf (Q)\texttt{M}-types and Coinduction in HoTT / CTT}\\[2ex]
  {\huge\sf Lasse Letager Hansen, 201912345}\\[2ex]
  \noindent\rule{\linewidth}{1mm}\\[4ex]
  \noindent{\Large\sf Master's Thesis, Computer Science\\[1ex] 
    \today \\[1ex] Advisor: Bas Spitters\\[15ex]}\\[\fill]}
\epsfig{file=logo.eps}\clearpage

%%%%%%%%%%%%%%%%%%%%%%%%%%%%%%%%%%%%%%%%%%%%%%%%%%%%%%%%%%%%%%%%%%%%%%% 

\pagestyle{plain}
\chapter*{Abstract}
\addcontentsline{toc}{chapter}{Abstract}

We present a construction of \texttt{M}-types from containers in a cubical type theory (CTT), We show how the containers construct a coalgebra, for which we can define a coinduction principle, making strong bisimulation imply equality. We then show constructions of \texttt{M}-types, and how they can be quotiented to construct what we call \texttt{QM}-types. The problem with \texttt{QM}-types is that in general assuming that we can lift function types of equivalence classes is equivalent to the axiom of choice \cite{DBLP:tt-in-tt}, but this can be solved by defining the quotienting relation and the type at the same time as a quotient inductive-inductive type (QIIT), which assuming the axiom of (countable) choice, is equal to the \texttt{QM}-type. We conclude with some examples of how to use \texttt{M}-types and some properties. All work is formalized in Cubical Agda, and the work on defining \texttt{M}-types has been accepted to the Cubical Agda github repository.

\chapter*{Resum\'e}
\addcontentsline{toc}{chapter}{Resum\'e}

\todo[inline]{in Danish\dots}

\chapter*{Acknowledgments}
\addcontentsline{toc}{chapter}{Acknowledgments}
I would like to thank my supervisor Bas Spitters for some enlightening discussions and quick responses and for taking the time to do weekly meetings regarding this masters thesis. I would also like to thank my fellow students, who made this challenging but fun time less lonely even in the time of the Corona pandemic.
\vspace{2ex}
\begin{flushright}
  \emph{Lasse Letager Hansen,}\\
  \emph{Aarhus, \today.}
\end{flushright}

\tableofcontents
\cleardoublepage
\pagenumbering{arabic}
\setcounter{secnumdepth}{3} % 2

%%%%%%%%%%%%%%%%%%%%%%%%%%%%%%%%%%%%%%%%%%%%%%%%%%%%%%%%%%%%%%%%%%%%%%% 

\chapter{Introduction} % Everyone interested should be able to read this!
\label{ch:intro}

\todo[inline]{motivate and explain the problem to be addressed}
\todo[inline]{get your bibtex entries from \url{https://dblp.org/}}

when \todo{...}... a constructor for a type \(\type{A}\) is a function, that takes some arguments, and returns an element \(a : \type{A}\), dually a destructor of \(\type{A}\) will return something given an element \(a : \type{A}\). \todo{...}... we can define a type inductively from a set of constructors, for example the natural numbers \(\mathbb{N}\), which are defined by the two constructors
\begin{center}
  \strut
  \hfill
  \begin{minipage}{0.15\linewidth}
    \begin{equation}
      \vcenter{\infer{0 : \mathbb{N}}{\strut}}
    \end{equation}
  \end{minipage}
  \hfill
  \begin{minipage}{0.20\linewidth}
    \begin{equation}
      \vcenter{\infer{\constructor{succ}~n : \mathbb{N}}{n : \mathbb{N}}}
    \end{equation}
  \end{minipage}
  \hfill
  \strut
\end{center}
with an equivalence relation \(\relation{\sim_{\mathbb{N}}}\) defined inductively, meaning we follow the structure of defined by the constructors
\begin{center}
  \strut
  \hfill
  \begin{minipage}{0.20\linewidth}
    \begin{equation}
      \vcenter{\infer[\constructor{\sim_{0}}]{0 = 0}{\strut}}
    \end{equation}
  \end{minipage}
  \hfill
  \begin{minipage}{0.35\linewidth}
    \begin{equation}
      \vcenter{\infer[\constructor{\sim_{succ}}]{\constructor{succ}~n~\relation{\sim_{\mathbb{N}}}~\constructor{succ}~m}{n~\relation{\sim_{\mathbb{N}}}~m}}
    \end{equation}
  \end{minipage}
  \hfill
  \strut
\end{center}
This relation gives equality, so if \(a~\relation{\sim_{\mathbb{N}}}~b\) then \(a \equiv b\). Likewise there is coinductive construction, where the focus is on the destructors instead of the constructors, an example is streams, representing an infinite sequence of elements, that can be defined from the two destructors head (\(\destructor{hd}\)) and tail (\(\destructor{tl}\)), where head represents the first element, and tail represents the rest of the sequence. The inference rules are given as
\begin{center}
  \strut
  \hfill
  \begin{minipage}{0.225\linewidth}
    \begin{equation}
      \vcenter{\infer{\destructor{hd}~\term{s} : \type{A}}{\term{s} : \typeformer{stream}~\type{A}}}
    \end{equation}
  \end{minipage}
  \hfill
  \begin{minipage}{0.25\linewidth}
    \begin{equation}
      \vcenter{\infer{\destructor{tl}~\term{s}: \typeformer{stream}~\type{A}}{\term{s} : \typeformer{stream}~\type{A}}}
    \end{equation}
  \end{minipage}
  \hfill
  \strut
\end{center}
we can again define an equivalence relation \(\relation{\sim_{stream}}\), but this time coinductively, focusing on the structure of the destructors instead
\begin{equation}
  \vcenter{\infer{\term{s}~\relation{\sim_{stream}}~\term{t}}{\destructor{hd}~\term{s} = \destructor{hd}~\term{t} & \destructor{tl}~\term{s}~\relation{\sim_{stream}}~\destructor{tl}~\term{t}}}
\end{equation}
This relation does not give an equality in MLTT, we just get bisimilarity meaning elements "behave" the same, but they are not equivalent. However one of the advantages of using homotopy type theory is that bisimilarity becomes equality. We will use this advantage to get the "correct" for of equality for some example of coinductive types. \todo{...}... An algebra is an operator \(\functor{F}\) with some closure relation \todo{cite something}. \todo{...}... inductive types are given as the initial algebra for some functor, this can be formalized as \texttt{W}-types, dually the coinductive types that we are interested in, can be formalized as the final coalgebra \todo{final coalgebra is scary words, find something more "noob" friendly} for some functor. We will be looking at how to define some coinductive types, as \texttt{M}-types, and define some bisimilarity relations for these types, showing we get equality when using homotopy type theory. We will then introduce weaker notion of bisimilarity, that does not yield equality, but can be used to construct a new type, by quotienting with the relation, giving us a type where the relation gives equality\todo{describe quotienting and its usefullness by a couple examples}. \todo{...}... When working with quotiented coinductive types. \todo{What is a monad?} \todo{...}... We define the quotiented delay monad \(\typeformer{Delay}/\relation{\sim_{weak}}\), and want to show that we can construct a partiality monad from this construction. A problem with the partiality operation \(D(-)/{\sim}_D\) is that countable choice is needed to show that it is a monad, however using QIIT types we can get around this problem, furthermore we can show that assuming countable choice, these two constructions are equal.

\subsection{A Brief Overview of the Thesis}
The rest of this chapter will introduce the some background theory and notation used in the rest of the thesis. In the second chapter we construct \texttt{M}-types from containers, and define a coinduction principle for the M-types. In the third chapter we give some example constructions of \texttt{M}-types. In the fourth chapter we introduce quotiented \texttt{M}-types (QM-types), and show equalities between these and quotient inductive-inductive types (QIITs), we show the construction of the partiality monad as an example. In the fifth and sixth chapter we go through various applications of the theory developed in the first couple chapters. Finally we conclude by discussion future research and improvements.

\section{Background Theory}
We start by giving some background theory / history on cubical type theory and summarize important concepts used in the rest of this thesis. \todo{Should a description of Set Theory be included}
\\ \\
We will be using \textbf{type theory} as the basis for mathematics. In type theory every term \(\term{x}\) is an element of some type \(\type{A}\), written \(\term{x} : \type{A}\). The idea in type theory is that propositions are types, so proofs boils down to showing that there exists an element of some type representing a proposition. Specifically proofs of equality becomes construction of an element of an equality type. The type theory we are working in is inspired by \textbf{Martin L\"of Type Theory (MLTT)} / Intuitionistic type Theory (ITT), which is designed on the principles of mathematical constructivism, where any existence proof must contain a witness. Meaning a proof of existence, can be converted into an algorithm that finds the element making the statement true. MLTT is built from the three finite types \(\empt\), \(\unit\) and \(\type{\mathbf{2}}\), and type constructors \(\Sigma\), \(\pi\) and \(=\). There is only a single way to make terms of \(=\)-type, and that is \(\mathtt{refl} : \prod_{a : A} (a = a)\). \todo{inductive types, universes and judgements} \todo{extensional vs intensional \url{https://ncatlab.org/nlab/show/intensional+type+theory}} \todo{find citation, currently \url{https://en.wikipedia.org/wiki/Intuitionistic_type_theory}}. We will be working in a cubical type theory, where the univalence axiom holds. The \textbf{univalence axiom} says that equality is equivalent to equivalence
\begin{equation}
  (\type{A} = \type{B}) \simeq (\type{A} \simeq \type{B})
\end{equation}
meaning if two objects are equivalent, then there is an equality between them, such that we can replace one by the other.
\\ \\
The type theory we will be working in is called \textbf{homotopy type theory (HoTT)} \cite{hottbook}, which is an intensional dependent type theory (built on MLTT) with the univalence axiom and higher inductive types. In HoTT the idenitity types form path spaces, so proofs of identity are not just \texttt{refl}, as for MLTT. Types are seen as "spaces", and we think of \(\term{a} : \type{A}\) as \(\term{a}\) being a point in the space \(\type{A}\), similarly functions are regarded as continuous maps from one space to another \cite{nlab:homotopy_type_theory}. The use of a type theory with the univalence axiom as a foundations for mathematics, is called \textbf{Univalent Foundations (UF)}\todo{should this be left out?}, which HoTT is. We will be working in a \textbf{cubical type theory (CTT)} \cite{DBLP:CTT} where the univalence axiom, is not just an axiom, but a statement that can actually be proven, meaning we can reduce the use of the univalence axiom, making it easier to do proofs involving the univalence axiom \cite{nlab:cubical_type_theory}. \todo{example of what computational "axioms" mean} The reason for the name cubical type theory, is because composition is defined by square, that is given three sides of a square we get the last one, see \figref{cubical-composition-square}.
\begin{figure}[h]
  \centering
  \begin{tikzcd}
    A \ar[r,"p \ct q \ct r",dashed] & B \\
    C \ar[u,"p^{-1}"] \ar[r,"q"'] & D \ar[u,"r"']
  \end{tikzcd}
  \caption{Composition square}
  \label{fig:cubical-composition-square}
\end{figure}
Most of the work in this thesis has been formalized in the proof assistant / programming languge Cubical Agda. A proof assistant helps with verifying proofs, while making the process of making proofs interactive. \textbf{Cubical Agda} is an implementation of a cubical type theory done by extending the proof assistant Agda. One of the main additions is the interval and path types. The interval can be thought of as elements in \([0,1]\). When working with the interval, we can only access the left and right endpoint \(\constant{i0}\) and \(\constant{i1}\) or some unspecified point in the middle \(i\), modeling the intuition of a continuous interval. Cubical agda also generalizes transporting, given a type line \(\typeformer{A} : \mathbb{I} \rightarrow \universe{\mathcal{U}}\), and the endpoint \(\typeformer{A}~\constant{i0}\) you get a line from \(\typeformer{A}~\constant{i0}\) to \(\typeformer{A}~\constant{i1}\) \cite{cubicalagda}. \todo{Axioms of cubical Agda} \todo{Some better (intuitive) description of the interval type!!}
\\ \\
Homotopy type theory is proof relevant, which means that a there might be multiple proofs of one statement, and these proofs might not be interchangeable (equal). The reason is that types in HoTT have a H-level, describing how equality behaves. We start from (-2) with contractiable types, meaning there is an element which all other element are equal to. Then there is (-1)-types which are mere propositions or hProp, where all elements of the type are equal, but there might not be any. If the type is inhabited, then we say the proposition is true. The 0-types are the hSets, where all equalities between two elements \(x,y\) are equal. For 1-types (1-groupoids) we get equalities of equalities are equal, and then so on for homotopy \(n\)-types. Any \(n\)-type is also a \(n+1\)-type, but with trivial equalities at the \(n+1\) level. If we don't want to do proof relevant mathematics we can do propositional truncation, converting types to \(-1\)-type, meaning we ignore the difference in proofs by just look at whether a type is inhabited or not. However doing this we loose some of the reasoning power of HoTT. One of the tools we get using the full power of HoTT is \textbf{Higher order inductive types (HITs)}, where we define types with point constructors and equality constructors, an example is the propositional truncation we just described\todo{write up constructors for propositional truncation}, another useful example is set truncated quotients\todo{write up constructors for set truncated quotients}.
\\ \\
A problem that arise when using a constructive type theory is the axiom of choice and the law of excluded middle, which does not have a computational interpretation, so to maintain the computational aspects of HoTT and CTT, we try to not use these axioms \cite[Introduction]{hottbook}.
\\ \\
\todo[inline]{W-types, Induction , M-types, Coinduction..., see introduction?}
Coinduction is the dual concept (in a categorical manner) of induction. The induction principle is an equivalence principle for congruent elements in an initial algebra.
\\ \\
Cubical Agda has a hierarchy of universes
\todo[inline]{Universes}

\subsubsection{UIP}
\paragraph{Path type}
We add a type \(\mathbb{I}\), which is defined to be a free de Morgan algebra on a discrete infinite set of names \todo{cite: https://arxiv.org/pdf/1611.02108.pdf}. The elements of \(\mathbb{I}\) can be described by the grammar
\begin{equation}
  r,s ::= 0 \mid 1 \mid i \mid 1 - r \mid r \land s \mid r \lor s
\end{equation}
The set \(\mathbb{I}\) has decidable equality. The elements in \(\mathbb{I}\) can be thought as formal representations of elements in the unit interval \([0,1]\). There is a special substitution with \(\constant{i0}\) and \(\constant{i1}\) being the endpoints of \([0,1]\).

\subsection{Problems with using AC and LEM}
These axioms does not have a computational interpretation, so to maintain the computational aspects of HoTT and CTT, we try to not use these axioms. \cite[Introduction]{hottbook}

\subsection{Universes ??}

\section{Notation}
Most of the work is formalized in the Cubical Agda proof assistant. The following is the notation / fonts used to denote specific definitions / concepts
\begin{itemize}
\item Universe \(\universe{\mathcal{U}}_i\) or \(\universe{\mathcal{U}}\)
\item Type \(\type{A} : \universe{\mathcal{U}}\)
\item A type former or dependent type \(\typeformer{B} : \type{A} \rightarrow \universe{\mathcal{U}}\)
\item A term \(\term{x} : \type{A}\) or for constants \(\constant{c} : \type{A}\)
\item A function \(\function{f} : \type{A} \rightarrow \type{C}\)
\item A constructor \(\constructor{f} : \type{A} \rightarrow \type{C}\) \todo{better description, not always a function}
\item A destructor \(\destructor{f} : \type{A} \rightarrow \type{C}\) \todo{better description, not always a function}
\item A path \(\pathterm{p} : \type{A} \equiv \type{C}\), heterogeneous paths are denotes \(\equiv_p\) or if the path is clear from context \(\equiv_*\).
\item A relation \(\relation{R} : \type{A} \rightarrow \type{A} \rightarrow \universe{\mathcal{U}}\) with notation \(\term{x}~\relation{R}~\term{y}\).
\item The unit type is \(\unit\) while the empty type is \(\empt\).
\item A functor \(\functor{P}\)
\item A container is denoted as \(\container{S}\) or \(\containerpair{\type{A}}{\typeformer{B}}\)
\item A coalgebra \(\coalg{\type{C}}{\function{\gamma}}\)
\item We denote the function giving the first and second projection of a dependent pair by \(\pi_1\) and \(\pi_2\).
\end{itemize}
Furthermore we define some useful notation for casing on natural numbers.
\begin{defn}
  \begin{equation}
    \natcases{\term{x}}{\function{f}} = \lambda\,n, \begin{cases} \term{x} & n = 0 \\ \function{f}~m & n = m + 1 \end{cases}
  \end{equation}
\end{defn}

%%%%%%%%%%%%%%%%%%%%%%%%%%%%%%%%%%%%%%%%%%%%%%%%%%%%%%%%%%%%%%%%%%%%%%% 

\chapter{\texttt{M}-types}
\label{ch:m-types}
In this chapter we will introduce containers (aka. signatures), and use them to construct \texttt{M}-types and operations \(\function{in}\) and \(\function{out}\) on the \texttt{M}-types (\thmref{M-shift}). We conclude the chapter by defining a coinduction principle for \texttt{M}-types \cite{DBLP:non-wellfounded}. \todo{Some formalization of coalgebras is missing?}

\section{Containers / Signatures}
\begin{defn}
  A Container (or signature) is a dependent pair \(\container{S} = \containerpair{\type{A}}{\typeformer{B}}\) for the types \(\type{A} : \universe{\mathcal{U}}\) and \(\typeformer{B} : \type{A} \rightarrow \universe{\mathcal{U}}\).
\end{defn}
\begin{defn}
  A polynomial functor over the container \(\container{S} = \containerpair{\type{A}}{\typeformer{B}}\) is defined for types as
  \begin{equation}
    \begin{aligned}
      &\functor{P}_{\container{S}} : \universe{\mathcal{U}} \rightarrow \universe{\mathcal{U}} \\
      &\functor{P}_{\container{S}}(\type{X}) = \sum_{\term{a} : \type{A}} \typeformer{B}(\term{a}) \rightarrow \type{X}
    \end{aligned}
  \end{equation}
  and for a function \(\function{f} : \type{X} \rightarrow \type{Y}\) as
  \begin{equation}
    \begin{aligned}
      &\functor{P}_{\container{S}}\function{f} : \functor{P}_{\container{S}}\type{X} \rightarrow \functor{P}_{\container{S}}\type{Y} \\
      &\functor{P}_{\container{S}}\function{f}(\term{a},\function{g}) = (\term{a}, \function{f} \circ \function{g}).
    \end{aligned}
  \end{equation}
\end{defn}
\begin{exmp}
  The polynomial functor for streams over the type \(\type{A}\) is defined by the container \(\container{S} = \containerpairsimple{\type{A}}{\unit}\), we get
  \begin{equation}
    \label{eq:stream-polyfunctor-pre}
    \functor{P}_{\container{S}}(\type{X}) = \sum_{\term{a} : \type{A}} \unit \rightarrow \type{X}.
  \end{equation}
  Since we are working in a logic with exponentials, we get \(\unit \rightarrow \type{X} \equiv \type{X}^{\unit} \equiv \type{X}\). Furthermore \(\unit\) and \(\type{X}\) does not depend on \(\type{A}\), so \eqref{eq:stream-polyfunctor-pre} is equivalent to
  \begin{equation}
    \functor{P}_{\container{S}}(\type{X}) = \type{A} \times \type{X}.
  \end{equation}
\end{exmp}
\noindent We now construct the \(\functor{P}_{\container{S}}\)-coalgebra for a polynomial functor \(\functor{P}_{\container{S}}\).
\begin{defn}
  A \(\functor{P}_{\container{S}}\)-coalgebra is defined as
  \begin{equation}
    \typeformer{Coalg}_{\container{S}} = \sum_{\type{C} : \universe{\mathcal{U}}} \type{C} \rightarrow \functor{P}_{\container{S}}\type{C}.
  \end{equation}
  We denote a \(\functor{P}_{\container{S}}\)-coalgebra given by \(\type{C}\) and \(\function{\gamma}\) as \(\coalg{\type{C}}{\function{\gamma}}\). Coalgebra morphisms are defined as
  \begin{equation}
    \begin{aligned}
      \cdot \Rightarrow ~\cdot &: \typeformer{Coalg}_{\container{S}} \rightarrow \typeformer{Coalg}_{\container{S}} \\
      \coalg{\type{C}}{\function{\gamma}} \Rightarrow \coalg{\type{D}}{\function{\delta}} &= \sum_{\function{f} : \type{C} \rightarrow \type{D}} \function{\delta} \circ \function{f} = \functor{P} \function{f} \circ \function{\gamma}
    \end{aligned}
  \end{equation}
\end{defn}
\noindent We can now define \texttt{M}-types.
\begin{defn}
  Given a container \(\container{S}\), we define \(\mathtt{M}\)-types as the type, making the coalgebra given by \(\typeformer{M}_{\container{S}}\) and \(\function{out} : \typeformer{M}_{\container{S}} \rightarrow \functor{P}_{\container{S}}(\typeformer{M}_{\container{S}})\) fulfill the property
  \begin{equation}
    \typeformer{Final}_{\container{S}} := \sum_{(\coalg{\type{X}}{\function{\rho}} : \mathtt{Coalg}_{\container{S}})} \prod_{(\coalg{\type{C}}{\function{\gamma}} : \mathtt{Coalg}_{\container{S}})} \mathtt{isContr}~(\coalg{\type{C}}{\function{\gamma}} \Rightarrow \coalg{\type{X}}{\function{\rho}}).
  \end{equation}
  That is \(\prod_{(\coalg{\type{C}}{\function{\gamma}} : \mathtt{Coalg}_{\container{S}})} \mathtt{isContr}(\coalg{\type{C}}{\function{\gamma}} \Rightarrow \coalg{\typeformer{M}_{\container{S}}}{\function{out}})\). We denote the \texttt{M}-type as \(\typeformer{M}_{\containerpair{\type{A}}{\typeformer{B}}}\) or \(\typeformer{M}_{\container{S}}\) or just \(\typeformer{M}\) when the Container is clear from the context.
\end{defn}
\noindent Continuing our example we now construct streams as an \(\mathtt{M}\)-type.
\begin{exmp} We define streams over the type \(\type{A}\) as the \(\mathtt{M}\)-type over the container \(\containerpairsimple{\type{A}}{\unit}\). If we apply the polynomial functor to the \texttt{M}-type, then we get \(\functor{P}_{\containerpairsimple{\type{A}}{\unit}} \typeformer{M} = \type{A} \times \typeformer{M}_{\containerpair{\type{A}}{\lambda\,\_,\unit}}\), illustrated in \figref{stream-M-type}.
  \begin{figure}[h]
    \centering
    \begin{tikzcd}[row sep=large,column sep=large]
      \type{A} & \type{A} \times \typeformer{M}_{\containerpair{A}{\lambda\,\_,\unit}} \ar[l,"\pi_1"'] \ar[r,"\pi_2"] & \typeformer{M}_{\containerpair{A}{\lambda\,\_,\unit}} \\
      & \typeformer{M}_{\containerpair{A}{\lambda\,\_,\unit}} \ar[ul,"\destructor{hd}"] \ar[u,"\function{out}"'] \ar[ur,"\destructor{tl}"']
    \end{tikzcd}
    \caption{\texttt{M}-types of streams}
    \label{fig:stream-M-type}
  \end{figure}
  We will that \(\function{out}\) is an isomorphism with inverse \(\function{in} : \functor{P}_{\container{S}}(\typeformer{M}) \rightarrow \typeformer{M}\) later in this section. We now have a semantic for the rules, we would expect for streams, if we let \(\constructor{cons} = \function{in}\) and \(\typeformer{stream}~\type{A} = \typeformer{M}_{\containerpair{\type{A}}{\lambda\,\_,\unit}}\),
  \begin{equation}
    \vcenter{\infer[\mathtt{E}_{\destructor{hd}}]{\destructor{hd}~\term{s} : \type{A}}{\type{A} : \universe{\mathcal{U}} & \term{s} : \typeformer{stream}~\type{A}}}
  \end{equation}
  \begin{equation}
    \vcenter{\infer[\mathtt{E}_{\destructor{tl}}]{\destructor{tl}~\term{s} : \typeformer{stream}~\type{A}}{\type{A} : \universe{\mathcal{U}} & \term{s} : \typeformer{stream}~\type{A}}}
  \end{equation}
  \begin{equation}
    \vcenter{\infer[\mathtt{I}_{\constructor{cons}}]{\term{\constructor{cons}}~\term{x}~\term{xs} : \typeformer{stream}~\type{A}}{\type{A} : \universe{\mathcal{U}} & \term{x} : \type{A} & \term{xs} : \typeformer{stream}~\type{A}}}
  \end{equation}
  or more precisely \(\destructor{hd} = \pi_1 \circ \function{out}\) and \(\destructor{tl} = \pi_2 \circ \function{out}\).
\end{exmp}
\begin{defn}
  We define a chain as a family of morphisms \(\function{\pi}_{(n)} : \type{X}_{n+1} \rightarrow \type{X}_{n}\), over a family of types \(\type{X}_{n}\). See \figref{x-chain}.
  \begin{figure}[h]
    \centering
    \begin{tikzcd}[row sep=large,column sep=large]
      \type{X}_0 & \type{X}_1 \ar[l,"\function{\pi}_{(0)}"] & \cdots \ar[l,"\function{\pi}_{(1)}"] & \type{X}_n \ar[l,"\function{\pi}_{(n-1)}"] & \type{X}_{n + 1} \ar[l,"\function{\pi}_{(n)}"] & \cdots \ar[l,"\function{\pi}_{(n+1)}"]
    \end{tikzcd}
    \caption{Chain of types / functions}
    \label{fig:x-chain}
  \end{figure}  
\end{defn}
\begin{lem}\label{lem:function-to-M-type-is-cone}
  For all coalgebras \(\coalg{\type{C}}{\function{\gamma}}\) for the container \(\container{S}\), we get \(\type{C} \rightarrow \typeformer{M}_{\container{S}} \equiv \typeformer{Cone}_{\coalg{\type{C}}{\function{\gamma}}}\), where \(\typeformer{Cone} = \sum_{(\function{f} : \prod_{(n : \mathbb{N})} \type{C} \rightarrow \type{X}_n)} \prod_{(n : \mathbb{N})} \function{\pi}_{(n)} \circ (\function{f}_{(n+1)}) \equiv f_n\) illustrated in \figref{M-seq-cone}.
  \begin{figure}[h]
    \centering
    \begin{tikzcd}[row sep=large,column sep=large]
      \type{X}_0 & \type{X}_1 \ar[l,"\function{\pi}_{(0)}"] & \cdots \ar[l,"\function{\pi}_{(1)}"] & \type{X}_n \ar[l,"\function{\pi}_{(n-1)}"] & \type{X}_{n + 1} \ar[l,"\function{\pi}_{(n)}"] & \cdots \ar[l,"\function{\pi}_{(n+1)}"] \\
      & & \type{C} \ar[ull,"\function{f}_0"] \ar[ul,"\function{f}_1"'] \ar[u,draw=none,"\cdots" description] \ar[ur,"\function{f}_n"] \ar[urr,"\function{f}_{n+1}"'] \ar[urrr,draw=none,"\cdots"'] & 
    \end{tikzcd}
    \caption{Cone}
    \label{fig:M-seq-cone}
  \end{figure}  
  \begin{proof}
    We define an isomorphism from \(\type{C} \rightarrow \typeformer{M}_{\container{S}}\) to \(\typeformer{Cone}_{\coalg{\type{C}}{\function{\gamma}}}\)
    \begin{align}
      &\function{fun}_{\pathterm{collapse}}~\function{f} = (\lambda\,n\,\term{z},\pi_1~(\function{f}~\term{z})~n),(\lambda\,n~i~\term{a},~\pi_2~(\function{f}~\term{a})~n~i) \\
      &\function{inv}_{\pathterm{collapse}}~(\function{u},\function{q})~\term{z} = (\lambda\,n,~\function{u}~n~\term{z}),(\lambda\,n\,i,~\function{q}~n~i~\term{z}) \\
      &\function{rinv}_{\pathterm{collapse}}~(\function{u},\function{q}) = \mathtt{refl}_{(\function{u},\function{q})} \\
      &\function{linv}_{\pathterm{collapse}}~\function{f} = \mathtt{refl}_{\function{f}}
    \end{align}
  \end{proof}
\end{lem}
\begin{lem}\label{lem:limit-collapse}
  Given \(\function{\ell} : \prod_{(n : \mathbb{N})} (\type{X}_{n} \rightarrow \type{X}_{n+1})\) and \(\term{y} : \sum_{\left(\term{x} : \prod_{(n : \mathbb{N})} \type{X}_n\right)} \term{x}_{n + 1} \equiv \function{\ell}_n~\term{x}_n\) the chain collapses as the equality \(\type{\mathcal{L}} \equiv \type{X}_0\).
  \begin{proof}
    We define this collapse by the isomorphism
    \begin{align}
      &\function{fun}_{\pathterm{\mathcal{L}collapse}}~(\term{x} , \pathterm{r}) = \term{x}_0 \\
      &\function{inv}_{\pathterm{\mathcal{L}collapse}}~\term{x}_0 = (\lambda\,n,~\function{\ell}^{(n)}~\term{x}_0)~,~(\lambda\,n,~\mathtt{refl}_{(\function{\ell}^{(n+1)}\,\term{x}_0)}) \\
      &\function{rinv}_{\pathterm{\mathcal{L}collapse}}~\term{x}_0 = \mathtt{refl}_{\term{x}_0}
    \end{align}
    where \(\function{\ell}^{(n)} = \function{\ell}_n \circ \function{\ell}_{n-1} \circ \dots \circ \function{\ell}_1 \circ \function{\ell}_0\). To define \(\function{linv}_{\pathterm{\mathcal{L}collapse}}~(\term{x} , \pathterm{r})\), we first define a fiber \((\type{X} , \term{z}, \function{\ell})\) over \(\mathbb{N}\) given some \(\term{z} : \type{X}_0\). Then any element of the type \(\sum_{\left(\term{x} : \prod_{(n : \mathbb{N})} \type{X}_n\right)} \term{x}_{n + 1} \equiv \function{\ell}_n~\term{x}_n\) is equal to a section over the fiber we defined. This means \(\term{y}\) is equal to a section. Since the sections are defined over \(\mathbb{N}\), which is an initial algebra for the functor \(\functor{G}\type{Y} = \unit + \type{Y}\), we get that sections are contractible, meaning \(\term{y} \equiv \function{inv}_{\mathcal{L}collapse} (\function{fun}_{\mathcal{L}collapse}~\term{y})\), since both are equal to sections over \(\mathbb{N}\).
\end{proof}
\end{lem}
\noindent We can now define the construction of \function{in} and \function{out}.
\begin{thm}
  \label{thm:M-shift}
   Given the container \(\containerpair{\type{A}}{\typeformer{B}}\) we define the equality
  \begin{equation}
    \pathterm{shift} : \type{\mathcal{L}} \equiv \functor{P}\type{\mathcal{L}}
  \end{equation}
  where \(\functor{P}\type{\mathcal{L}}\) is the limit of a shifted sequence. Then
  \begin{equation}
    \function{in} = \mathtt{transport}~\pathterm{shift}
  \end{equation}
  \begin{equation}
    \function{out} = \mathtt{transport}~(\sym{\pathterm{shift}}).
  \end{equation}
  \begin{proof}
    The proof is done using the two helper lemmas
    \begin{equation}
      \pathterm{\alpha} : \type{\mathcal{L}}^{\functor{P}} \equiv \functor{P}\type{\mathcal{L}}
    \end{equation}
    \begin{equation}
      \pathterm{\mathcal{L}unique} : \type{\mathcal{L}} \equiv \type{\mathcal{L}}^{\functor{P}}
    \end{equation}
    We define \(\pathterm{\mathcal{L}unique}\) by the ismorphism
    \begin{align}
      \function{fun}_{\pathterm{\mathcal{L}unique}}~(\function{a} , \function{b}) &= \natcases{\unitelem}{\function{a}} , \natcases{\mathtt{refl}_{\unitelem}}{\function{b}} \\
    \function{inv}_{\pathterm{\mathcal{L}unique}}~(\function{a} , \function{b}) &= \function{a} \circ \constructor{succ}~,~\function{b} \circ \constructor{succ} \\
    \function{rinv}_{\pathterm{\mathcal{L}unique}}~(\function{a} , \function{b}) &= \mathtt{refl}_{(\function{a},\function{b})} \\
    \function{linv}_{\pathterm{\mathcal{L}unique}}~(\function{a} , \function{b}) &= \mathtt{refl}_{(\function{a},\function{b})}
    \end{align}
    The definition of \(\pathterm{\alpha}\) is then, 
    \begin{align}
      \type{\mathcal{L}}^{\functor{P}} &\equiv \sum_{(\term{x} : \prod_{(n : \mathbb{N})} \sum_{(\term{a} : \type{A})} \typeformer{B}\,\term{a} \rightarrow \type{X}_n)} \prod_{(n : \mathbb{N})} \function{\pi}_{(n+1)}~\term{x}_{n+1} \equiv \term{x}_{n} \\
      &\equiv \sum_{(\term{x} : \sum_{(\term{a} : \prod_{(n : \mathbb{N})} \type{A})} \prod_{(n : \mathbb{N})} \term{a}_{n+1} \equiv \term{a}_n)} \sum_{(\function{u} : \prod_{(n : \mathbb{N})} \typeformer{B}\,{(\pi_1\,\term{x})}_n \rightarrow \type{X}_n)} \prod_{(n : \mathbb{N})} \function{\pi}_{(n)} \circ \function{u}_{n+1} \equiv_* \function{u}_n \label{eq:alpha-unfold} \\
      &\equiv \sum_{(\term{a} : \type{A})} \sum_{(\function{u} : \prod_{(n : \mathbb{N})} \typeformer{B}\,a \rightarrow \type{X}_n)} \prod_{(n : \mathbb{N})} \function{\pi}_{(n)} \circ \function{u}_{n+1} \equiv \function{u}_n \label{eq:alpha-limit-collapse} \\
      &\equiv \sum_{(\term{a} : \type{A})} \typeformer{B}~\term{a} \rightarrow \type{\mathcal{L}} \label{eq:alpha-cone}\\
      &\equiv \functor{P}\type{\mathcal{L}}
    \end{align}
    To collapse \(\sum_{(\term{a} : \prod_{(n : \mathbb{N})} \type{A})} \prod_{(n : \mathbb{N})} \term{a}_{n+1} \equiv \term{a}_n\) to \(\type{A}\) between \eqref{eq:alpha-unfold} and \eqref{eq:alpha-limit-collapse} we use \lemref{limit-collapse} . We use \lemref{function-to-M-type-is-cone} for the equality between \eqref{eq:alpha-limit-collapse} and \eqref{eq:alpha-cone}. The rest of the equalities are trivial. The definition of \(\pathterm{shift}\) is
    \begin{equation}
      \pathterm{shift} = \sym{\pathterm{\alpha}} \ct \pathterm{\mathcal{L}unique}.
    \end{equation}
    We furthermore get the definitions \(\function{in} = \mathtt{transport}~\pathterm{shift}\) and \(\function{out} = \mathtt{transport}~(\sym{\pathterm{shift}})\), since \function{in} and \function{out} are part of an equality relation \(\pathterm{shift}\), they are both surjective and embeddings. \todo{is surjectivity and embedding important here? Describe this where relevant instead!}
  \end{proof}
\end{thm}

\section{Coinduction Principle for \texttt{M}-types}
We can now construct a coinduction principle given a bisimulation relation.
\begin{defn}
  For all coalgebras \(\coalg{\type{C}}{\function{\gamma}} : \typeformer{Coalg}_{\container{S}}\), given a relation \(\relation{\mathcal{R}} : \type{C} \rightarrow \type{C} \rightarrow \universe{\mathcal{U}}\) and a type \(\type{\overline{\mathcal{R}}} = \sum_{(\term{a} : \type{C})} \sum_{(\term{b} : \type{C})} \term{a}~\relation{\mathcal{R}}~\term{b}\), such that \(\type{\overline{\mathcal{R}}}\) and \(\function{\alpha}_{\relation{\mathcal{R}}} : \type{\overline{\mathcal{R}}} \rightarrow \functor{P}_{\container{S}}(\type{\overline{\mathcal{R}}})\) forms a \(\functor{P}\)-coalgebra \(\coalg{\type{\overline{\mathcal{R}}}}{\function{\alpha}_{\relation{\mathcal{R}}}} : \mathtt{Coalg}_{\container{S}}\), making the diagram in \figref{coalgebra-coinduction-definition} commute \todo{What does commute mean here?} (\(\Longrightarrow\) represents \(\functor{P}\)-coalgebra morphisms).
  \begin{figure}[h]
    \centering
    \begin{tikzcd}[row sep=huge,column sep=huge]
      \coalg{\type{C}}{\function{\gamma}} & \coalg{\type{\overline{\mathcal{R}}}}{\function{\alpha}_{\relation{\mathcal{R}}}} \ar[l,Rightarrow,"{\pi_1}^{\type{\overline{\mathcal{R}}}}"'] \ar[r,Rightarrow,"{\pi_2}^{\type{\overline{\mathcal{R}}}}"] & \coalg{\type{C}}{\function{\gamma}}
    \end{tikzcd}
    \caption{Bisimulation for a coalgebra}
    \label{fig:coalgebra-coinduction-definition}
  \end{figure}
\end{defn}
\begin{defn}[Coinduction principle]
  \noindent Given a relation \(\relation{\mathcal{R}}\), that is part of a bisimulation over a final \(\functor{P}\)-coalgebra \(\coalg{\typeformer{M}}{\function{out}}:\mathtt{Coalg}_{\container{S}}\) we get the diagram in \figref{final-coalgebra-coinduction},
  \begin{figure}[h]
    \centering
    \begin{tikzcd}[row sep=huge,column sep=huge]
      \coalg{\typeformer{M}}{\function{out}} & \coalg{\type{\overline{\mathcal{R}}}}{\function{\alpha}_{\relation{\mathcal{R}}}} \ar[l,Rightarrow,"{\pi_1}^{\type{\overline{\mathcal{R}}}}"'] \ar[r,Rightarrow,"{\pi_2}^{\type{\overline{\mathcal{R}}}}"] & \coalg{\typeformer{M}}{\function{out}}
    \end{tikzcd}
    \caption{Bisimulation principle for final coalgebra}
    \label{fig:final-coalgebra-coinduction}
  \end{figure}
  \noindent where \(\pi_1^{\type{\overline{\mathcal{R}}}} = ~! = \pi_2^{\type{\overline{\mathcal{R}}}}\) \todo{what is \(!\)} where \(!\) is the unique mapping property (UMP) out of the final coalgebra. Given \(\term{r} : \term{m}~\relation{\mathcal{R}}~\term{m'}\) we get the equation
  \begin{equation}
    \term{m} = \pi_1^{\type{\overline{\mathcal{R}}}} (\term{m}, \term{m'}, \term{r}) = \pi_2^{\type{\overline{\mathcal{R}}}} (\term{m}, \term{m'}, \term{r}) = \term{m'}.
  \end{equation}
  \todo{What is the consequence of this?}
\end{defn}

%%%%%%%%%%%%%%%%%%%%%%%%%%%%%%%%%%%%%%%%%%%%%%%%%%%%%%%%%%%%%%%%%%%%%%% 

\chapter{Instantiation of M-types}
In this section we show some examples of types that can be constructed as \(M\)-types, and show how their constructors can be defined. We then conclude the chapter with some general observation, and define some rules for how to construct \texttt{M}-types. \todo{Is there anything else that is show for each M-type?}

\todo{complete this section}

\section{Stream Formalization using \texttt{M}-types}
As described earlier, given a type \(\type{A}\) we define the stream of that type as
\begin{equation}
  \typeformer{stream}~\type{A} := \typeformer{M}_{\containerpair{\type{A}}{\lambda\,\_,\unit}}
\end{equation}
this is equal to an alternative definition of streams \todo{...}

\section{ITrees as \texttt{M}-types}
Interaction trees (ITrees) \cite{DBLP:itrees} are used to model effectful behavior, where computations can interact with an external environment by events. ITrees are defined by the following constructors
\begin{equation}
  \vcenter{\infer[\mathtt{I}_{\constructor{Ret}}]{\constructor{Ret}~\term{r} : \typeformer{itree}~\typeformer{E}~\type{R}}{r : \type{R}}}
\end{equation}
\begin{equation}
  \vcenter{\infer[\mathtt{I}_{\constructor{Vis}}]{\constructor{Vis}~\term{a}~\term{f} : \typeformer{itree}~\typeformer{E}~\type{R}}{\type{A} : \universe{\mathcal{U}} & \term{a} : \typeformer{E}~\type{A} & \term{f} : \type{A} \rightarrow \typeformer{itree}~\typeformer{E}~\type{R}}}.
\end{equation}
\begin{equation}
  \vcenter{\infer[\mathtt{E}_{\constructor{Tau}}]{\constructor{Tau}~\term{t} : \typeformer{itree}~\typeformer{E}~\type{R}}{\term{t} : \typeformer{itree}~\typeformer{E}~\type{R}}}.
\end{equation}
where \(\type{R}\) is the type for returned values, while \(\typeformer{E}\) is a dependent type for events representing external interactions.

\subsection{Delay Monad}
We start by looking at ITrees without the \(\constructor{Vis}\) constructor, this type is also know as the delay monad. It can be used to model delayed computations, either returning immediately given by the constructor \(\constructor{now} = \constructor{Ret}\), or delayed some (possibly infinite) number of steps by the constructor \(\constructor{later} = \constructor{Tau}\). We construct this type as an \texttt{M}-type.
\begin{defn}
  The  delay monad can be defined as the \texttt{M}-type for the container
  \begin{equation}
    \container{S} = \containerpair{\type{R} +\unit}{\begin{cases} \empt & \mathtt{inl}~\term{r} \\ \unit & \mathtt{inr}~\unitelem \end{cases}}
  \end{equation}
\end{defn}
\noindent The polynomial functor for this container is
\begin{equation}
  \functor{P}_{\container{S}}(\type{X}) = \sum_{(\term{x} : \type{R} + \unit)} \begin{cases} \empt & \term{x} = \mathtt{inl}~\term{r} \\ \unit & \term{x} = \mathtt{inr}~\unitelem \end{cases} \rightarrow \type{X},
\end{equation}
which we can simplify
\begin{equation}
  \functor{P}_{\container{S}}(\type{X}) = \type{R} \times (\empt \rightarrow \type{X}) + \type{X}.
\end{equation}
We know that \((\empt \rightarrow \type{X}) \equiv \unit\), so we can simplify further to
\begin{equation}
  \functor{P}_{\container{S}}(\type{X}) = \type{X} + \type{R}
\end{equation}
meaning we get diagram in \figref{delay-monad}.
\begin{figure}[h]
  \centering
  \begin{tikzcd}[row sep=large,column sep=large]
    \type{R} \ar[dr,"\constructor{now}"'] \ar[r,"\texttt{inl}"] & \type{R} + \typeformer{M} \ar[d,"\function{in}"'] & \typeformer{M} \ar[l,"\texttt{inr}"'] \ar[dl,"\constructor{later}"] \\
    & \typeformer{M}
  \end{tikzcd}
  \caption{Delay monad}
  \label{fig:delay-monad}
\end{figure}
\noindent We can define the constructors \(\constructor{now}\) and \(\constructor{later}\) using \(\function{in}\) function for \texttt{M}-types, together with the injections \(\mathtt{inl}\) and \(\mathtt{inr}\).

\subsection{Tree}
Now lets look at the example, where we remove the \(\constructor{Tau}\) constructor. This gives us a type of tree, with leaves given by \(\constructor{Ret}\), and nodes given by \(\constructor{Vis}\) branching based on some type \(\type{A}\), for an event \(\term{a} : \typeformer{E}~\type{A}\). 
\begin{defn}
  We can define \(\type{R}\)-valued \(\typeformer{E}\)-event trees as the \texttt{M}-type defined by the container
  \begin{equation}
    \container{S} = \containerpair{\type{R} + \sum_{(\type{A} : \universe{\mathcal{U}})} \typeformer{E}~\type{A}}{\begin{cases}\empt & \mathtt{inl}~\term{r} \\ \type{A} & \mathtt{inr}~(\type{A},\term{e})\end{cases}}.
  \end{equation}
\end{defn}
\noindent The polynomial functor for this container is
\begin{equation}
  \functor{P}_{\container{S}}(\type{X}) = \sum_{(\term{x} : \type{R} + \sum_{(\type{A} : \universe{\mathcal{U}})} \typeformer{E}\,\type{A})} \begin{cases}\empt & \term{x} = \mathtt{inl}~\term{r} \\ \type{A} & \term{x} = \mathtt{inr}~(\type{A},\term{e})\end{cases} \rightarrow \type{X},
\end{equation}
which simplifies to
\begin{equation}
  \functor{P}_{\container{S}}(\type{X}) = (\type{R} \times (\empt \rightarrow \type{X})) + (\sum_{\type{A} : \universe{\mathcal{U}}} \typeformer{E}~\type{A} \times (\type{A} \rightarrow \type{X})),
\end{equation}
and further
\begin{equation}
  \functor{P}_{\container{S}}(\type{X}) = \type{R} + \sum_{\type{A} : \universe{\mathcal{U}}} \typeformer{E}~\type{A} \times (\type{A} \rightarrow \type{X}).
\end{equation}
We get the diagram in \figref{tree-constructors} for the \(\functor{P}\)-coalgebra.
\begin{figure}[h]
  \centering
  \begin{tikzcd}[row sep=large,column sep=large]
    \type{R} \ar[dr,"\constructor{Ret}"'] \ar[r,"\texttt{inl}"] & \type{R} + \sum_{\type{A} : \universe{U}} \typeformer{E}~\type{A} \times (\type{A} \rightarrow \typeformer{M}) \ar[d,"\function{in}"]  & \sum_{\type{A} : \universe{U}} \typeformer{E}~\type{A} \times (\type{A} \rightarrow \typeformer{M}) \ar[l,"\texttt{inr}"'] \ar[dl,"\constructor{Vis}"] \\
    & \typeformer{M}
  \end{tikzcd}
  \caption{Tree Constructors}
  \label{fig:tree-constructors}
\end{figure}
\noindent Again we can define \(\constructor{Ret}\) and \(\constructor{Vis}\) using the \(\function{in}\) function.

\subsection{ITrees}
\todo[inline]{Get the correct equivalence for ITrees (Part of project description?)}
Now we should have all the knowledge needed to make ITrees using \texttt{M}-types.
\begin{defn}
  We define the type of ITrees as the \texttt{M}-type given by the container
  \begin{equation}
    \container{S} = \containerpair{\type{R}+\unit+\sum_{\type{A}:\universe{\mathcal{U}}} (\typeformer{E}~\type{A})}{\begin{cases} \empt & \texttt{inl}~r \\ \unit & \texttt{inl}~(\texttt{inl}~\unitelem) \\ \type{A} & \texttt{inr}~(\texttt{inr}~(\type{A},\term{e})) \end{cases}}.
  \end{equation}
\end{defn}
\noindent The (reduced) polynomial functor for this container is
\begin{equation}
  \functor{P}_{\container{S}}(\type{X}) = \type{R} + \type{X} + \sum_{(\type{A} : \universe{\mathcal{U}})} (\typeformer{E}~\type{A} \times (\type{A} \rightarrow \type{X}))
\end{equation}
Giving us the diagram in \figref{ITree-constructors}, from which the constructors of the type can be defined using \(\function{in}\).
\begin{figure}[h]
  \centering
  \begin{tikzcd}[row sep=large,column sep=large]
    \type{R} \ar[ddrr,"\constructor{Ret}"'] \ar[r,"\texttt{inl}\,\circ\,\texttt{inr}"] & \typeformer{M} + \type{R} + \sum_{(\type{A} : \universe{\mathcal{U}})} (\typeformer{E}~\type{A} \times (\type{A} \rightarrow \typeformer{M})) \ar[ddr,"\function{in}"] & & \sum_{\type{A} : \universe{U}} \typeformer{E}~\type{A} \times (\type{A} \rightarrow \typeformer{M}) \ar[ll,"\texttt{inr}"'] \ar[ddl,"\constructor{Vis}"] \\
    & & \typeformer{M} \ar[ul,"\texttt{inl}\,\circ\,\texttt{inl}"'] \ar[d,"\constructor{Tau}"'] \\
    & & \typeformer{M}
  \end{tikzcd}
  \caption{ITree constructors}
  \label{fig:ITree-constructors}
\end{figure}

\section{Automaton}
An automaton is defined as a set of state \(\type{V}\) and an alphabet \(\type{\alpha}\) and a transition function \(\function{\delta} : \type{V} \rightarrow \type{\alpha} \rightarrow \type{V}\). This gives us the diagram in \figref{automaton}
\begin{figure}[h]
  \centering
  \begin{tikzcd}[row sep=huge,column sep=huge]
    \type{V} \times (\type{V} \rightarrow \type{\alpha} \rightarrow \type{V}) \ar[dr,"build"'] \ar[r] & M & M \ar[dl,"step"] \ar[l] \\
    & M \ar[u]
  \end{tikzcd}
  \caption{automaton}
  \label{fig:automaton}
\end{figure}

\section{General rules for constructing \texttt{M}-types}
We want to create a calculus for defining coinductive types as \texttt{M}-types. We would like to be able to define a type that has a given set of constructor / destructors rules. If we for example is given the rule
\begin{equation}
  \vcenter{\infer{\constructor{ret}~a : \type{T}}{\term{a} : \type{A}}}
\end{equation}
we get that it corresponds to the M-type for the container \(\containerpairsimple{\type{A}}{\empt}\), while if we have something that produces an element of it self as
\begin{equation}
  \vcenter{\infer{\destructor{tl}~\term{a} : \type{T}}{\term{a} : \type{T}}}
\end{equation}
the container is \(\containerpairsimple{\unit}{\unit}\) \todo{why give some arguments or illustation showing these reasonings}. If we want a type with both these rules, then we just take the disjoint union of the two containers
\begin{equation}
  \containerpair{\type{A} + \unit}{\begin{cases} \empt & \mathtt{inl}~\term{a} \\ \unit & \mathtt{inr}~\unitelem \end{cases}}
\end{equation}
which is the delay type. We can also define some more involved constructors, that build on the type itself
\begin{equation}
  \vcenter{\infer{\constructor{node}~\term{a} : \type{T}}{\term{a} : \type{A} \rightarrow \type{T}}}
\end{equation}
has container \(\containerpair{\unit}{\type{A}}\). We can types with a given destructor
\begin{equation}
  \vcenter{\infer{\destructor{hd}~\term{a} : \type{A}}{\term{a} : \type{T}}}
\end{equation}
has container \(\containerpairsimple{\type{A}}{\empt}\), but this is the same as for \(\constructor{ret}\), and we do not always want both. The difference is how they are added to other constructors / destructors. Destructors are easily added together, take for example \(\destructor{hd}\) and \(\destructor{tl}\), \(\containerpairsimple{\type{A}}{\unit}\). In general adding containers \(\containerpair{\type{A}}{\typeformer{B}}\) and \(\containerpair{\type{C}}{\typeformer{D}}\) for two constructors together is done by
\begin{equation}
  \containerpair{\type{A} + \type{C}}{ \begin{cases} \typeformer{B}~\term{a} & \mathtt{inl}~\term{a} \\ \typeformer{D}~\term{c} & \mathtt{inr}~\term{c} \end{cases}}
\end{equation}
whereas adding containers for two destructors is done by \todo{Is this correct? Seems correct..}
\begin{equation}
  \containerpairsimple{\type{A} \times \type{C}}{\lambda\,(\term{a},\term{c}),~\typeformer{B}~\term{a} + \typeformer{D}~\term{c}}  
\end{equation}
however combining destructors and constructors is not as simple. Anything type \(\type{T}\) that is defined using a record (except higher inductive types), will also be definable as an \texttt{M}-type. Given a record, which is a list of fields \(\term{f}_1 : \type{F}_1,\term{f}_2 : \type{F}_2,\dots,\term{f}_n : \type{F}_n\), we can construct the \texttt{M}-types by the container
\begin{equation}
  \containerpairsimple{\type{F}_1 \times \type{F}_2 \times \dots \times \type{F}_n}{\empt}
\end{equation}
where each destructor \(\destructor{d_n} : \type{T} \rightarrow \type{F_n}\) for the field \(\term{f}_n\) will be defined as \(\destructor{d_n}~\term{t} = \pi_n~(\function{out}~\term{t})\). However fields in a coinductive container may depend on previous defined fields, as given by the general list of fields \(\term{f}_1 : \type{F}_1,\term{f}_2 : \typeformer{F}_2,\cdots,\term{f}_n : \typeformer{F}_n\), where each field depends on all the previous once, this can be defined by the container
\begin{equation}
  \containerpairsimple{\sum_{(\term{f}_1 : \type{F}_1)} \sum_{(\term{f_2} : \typeformer{F}_2)} \cdots \sum_{(\term{f}_{n-1} : \typeformer{F}_{n-1})}\typeformer{F}_n}{\empt}
\end{equation}
however, if any of the destructors/fields are non dependent, then the can be added as a product (\(\times\)) instead of a dependent product (\(\Sigma\)). Furthermore the fields may construct an element of the type of the record \(\type{T}\), however anything after that field cannot on it, since it will break the strictness requirements of the record / coinductive type. As an example let \(\type{f}_1\) be a type and \(\function{f_2}\) be the function with type \(\typeformer{F}_2 = \type{f}_1 \rightarrow (\type{f}_1 \rightarrow \type{A}) \rightarrow \mathtt{M}\), which by currying is equal to \(\type{f}_1 \times (\type{f}_1 \rightarrow \type{A}) \rightarrow \mathtt{M}\), we can then define by the container
\begin{equation}
  \containerpair{\sum_{(\type{f}_1 : \universe{\mathcal{U}})} \left(\unit \times \sum_{(\term{f}_3 : \typeformer{F}_3)} \cdots \sum_{(\term{f}_{n-1} : \typeformer{F}_{n-1})}\typeformer{F}_n\right)}{\lambda\,(\type{f}_1,\unitelem,\term{f}_3,\dots),\typeformer{F}_2}
\end{equation}
where \(\typeformer{F}_2\) have been moved to the last part of the container, we can even leave out the "\(\unit \times\)" from the container. The types of the field can also be dependent \(\typeformer{F_2} = (\term{x} : \type{f}_1) \rightarrow \typeformer{B}~\term{x} \rightarrow \mathtt{M}\), but again by currying we can get \(\typeformer{F_2} : \sum_{(\term{x} : \type{f}_1)} \typeformer{B}~\term{x} \rightarrow \mathtt{M}\) which is defined by the container
\begin{equation}
  \containerpair{\sum_{(\type{f}_1 : \universe{\mathcal{U}})} \sum_{(\term{f}_3 : \typeformer{F_3})} \cdots \sum_{(\term{f}_{n-1} : \typeformer{F_{n-1}})}\typeformer{F_n}}{\lambda\,(\type{f}_1,\term{f}_3,\dots),\sum_{\term{x} : \type{f}_1} (\typeformer{B}~\term{x})}
\end{equation}
so we would also expect that a type defined as a (coinductive) record is equal to the version defined as a \texttt{M}-type. \todo{which is the case, proof needed!}
\\ \\
But we run into problems if ... \todo{Problem cases} \todo{What differes from W and M types for closure of constructors / destructors?}

\section{Wacky \texttt{M}-type}
We end this chapter by showing of some wacky \texttt{M}-type, that utilizes the definition of the \texttt{M}-type to the fullest.
\begin{defn}
  We define a wacky \texttt{M}-type by the following container
  \begin{equation}
    \containerpair{\mathbb{N} + \mathbb{N}}{\begin{cases} \mathbb{N} & \mathtt{inl}~0 \\ \empt & \mathtt{inl}~\term{x} \land \term{x}~\text{is odd} \\ \empt & \mathtt{inr}~\term{x} \land \term{x}~\text{is even} \\ \unit & o.w. \end{cases}}
  \end{equation}
\end{defn}
this container gives the \texttt{M}-type and constructors / destructors shown in \figref{wacky-m-type}
\begin{figure}[h]
  \centering
  \begin{tikzcd}[row sep=huge,column sep=huge]
    \mathbb{N} \ar[dr,"\constructor{odd}"',very near start] \ar[r,"\mathtt{inl} (2n+1)"] & M & \mathbb{N} \times M \ar[dl,"\destructor{step}",very near start] \ar[l,"o.w."] \\
    \mathbb{N} \ar[r,"\constructor{even}"'] \ar[ur,"\mathtt{inr}(2n)",very near start] & M \ar[u] & (\mathbb{N} \rightarrow M) \ar[l,"\constructor{inl\text{-}zero}"] \ar[ul,"\mathtt{inl}~0 \times f"',very near start]
  \end{tikzcd}
  \caption{wacky \texttt{M}-type}
  \label{fig:wacky-m-type}
\end{figure}
The type can be interpreted as a stream of coproducts of natural numbers, that terminates whenever it is the left injection and even, or the right injection and odd, and whenever the left injection is zero, it splits in a branches indexed by the natural numbers.
\\ \\
This show that we can define coinductive types, with rather complex structure.
\todo[inline]{define a wacky M-type showing some of the use of M-types / complexity}


%%%%%%%%%%%%%%%%%%%%%%%%%%%%%%%%%%%%%%%%%%%%%%%%%%%%%%%%%%%%%%%%%%%%%%% 

\chapter{QM-types}
In this chapter we will introduce quotients, and show how to construct quotiented \texttt{M}-types which we call Q\texttt{M}-types. We show you can be constructed a QIIT that is equal to the Q\texttt{M}-type assuming axiom of choice. \todo{better introduction to chapter, and reference to main points}

\section{Quotienting and Constructors}
\todo[inline]{Describe set truncated quotients and their construction / elimination principles, and how it relates to quotienting M-types}
\section{Quotient \texttt{M}-type}
Since we know that \texttt{M}-types preserves the H-level, we can use set-truncated quotients, to define quotient \texttt{M}-types, though we run into the problem of \todo{problem of direct quotients}, as a solution we use QIIT to define the relation and type at the same type.

\section{Quotient inductive-inductive types (QIITs)}
A quotient inductive-inductive type (QIIT) is a type together with a relation defined on that type. QIITs are HIITs that are set truncated.

\section{Partiality monad}
In this section we will define the partiality monad (see below) and show that (assuming the axiom of countable choice) the delay monad quotiented by weak bisimilarity.
\begin{defn}[Partiality Monad]
  A simple example of a quotient inductive-inductive type is the partiality monad \((-)_\bot\) over a type \(\type{R}\), defined by the constructors\\[-9mm]
  \begin{center}
    \begin{minipage}{0.25\linewidth}
      \begin{equation}
        \vcenter{\infer{\type{R}_\bot : \universe{\mathcal{U}}}{\strut}}
      \end{equation}
    \end{minipage}
    \hfill
    \begin{minipage}{0.25\linewidth}
      \begin{equation}
        \vcenter{\infer{\constructor{\bot} : \type{R}_\bot}{\strut}}
      \end{equation}
    \end{minipage}
    \hfill
    \begin{minipage}{0.25\linewidth}
      \begin{equation}
        \vcenter{\infer{\constructor{\eta}~\term{a} : \type{R}_\bot}{\term{a} : \type{R}}}
      \end{equation}
    \end{minipage}
  \end{center}
  and a relation \((\cdot~\relation{\sqsubseteq_{\bot}}~\cdot)\) indexed twice over \(\type{R}_\bot\), with properties\\[-9mm] \todo{Should I define what it means to be an ordering relation separately, and just say the relation here is an instance of that? (Generalize?)}
  \begin{center}
    \begin{minipage}{0.45\linewidth}
      \begin{equation}
        \vcenter{\infer{\constructor{\bigsqcup}~(\function{s},\function{b}) : \type{R}_\bot}{\function{s} : \mathbb{N} \rightarrow \type{R}_\bot & \function{b} : \prod_{(n : \mathbb{N})} \function{s}_n~\relation{\sqsubseteq_{\bot}}~\function{s}_{n+1} }}
      \end{equation}
    \end{minipage}
    \hfill
    \begin{minipage}{0.45\linewidth}
      \begin{equation}
        \vcenter{\infer{\constructor{\alpha_{\bot}}~\term{p}~\term{q} : \term{x} \equiv \term{y}}{\term{x},\term{y} : \type{R}_\bot & \term{p} : \term{x}~\relation{\sqsubseteq_{\bot}}~\term{y} & \term{q} : \term{y}~\relation{\sqsubseteq_{\bot}}~\term{x}}}
      \end{equation}
    \end{minipage}
  \end{center}
  \strut\\[-15mm]
  \begin{center}
    \begin{minipage}{0.25\linewidth}
      \begin{equation}
        \vcenter{\infer[\constructor{\sqsubseteq_{refl}}]{\term{x}~\relation{\sqsubseteq_{\bot}}~\term{x}}{\term{x} : \type{R}_\bot}}
      \end{equation}
    \end{minipage}
    \hfill
    \begin{minipage}{0.35\linewidth}
      \begin{equation}
        \vcenter{\infer[\constructor{\sqsubseteq_{trans}}]{\term{x}~\relation{\sqsubseteq_{\bot}}~\term{z}}{\term{x}~\relation{\sqsubseteq_{\bot}}~\term{y} & \term{y}~\relation{\sqsubseteq_{\bot}}~\term{z}}}
      \end{equation}
    \end{minipage}
    \hfill
    \begin{minipage}{0.25\linewidth}
      \begin{equation}
        \vcenter{\infer[\constructor{\sqsubseteq_{never}}]{\constructor{\bot}~\relation{\sqsubseteq_{\bot}}~\term{x}}{\term{x} : \type{R}_\bot}}
      \end{equation}
    \end{minipage}
  \end{center}
  \strut\\[-15mm]
  \begin{center}
    \hfill
    \begin{minipage}{0.50\linewidth}
      \begin{equation}
        \vcenter{\infer{\prod_{(n : \mathbb{N})} \function{s}_n~\relation{\sqsubseteq_{\bot}}~\bigsqcup(\function{s},\function{b})}{\function{s} : \mathbb{N} \rightarrow \type{R}_\bot & \function{b} : \prod_{(n : \mathbb{N})} \function{s}_n~\relation{\sqsubseteq_{\bot}}~\function{s}_{n+1} }}
      \end{equation}
    \end{minipage}
    \hfill
    \begin{minipage}{0.3\linewidth}
      \begin{equation}
        \vcenter{\infer{\bigsqcup(\function{s},\function{b})~\relation{\sqsubseteq_{\bot}}~\term{x}}{\prod_{(n:\mathbb{N})} \function{s}_n~\relation{\sqsubseteq_{\bot}}~\term{x}}}
      \end{equation}
    \end{minipage}
  \end{center}
  and finally set truncated
  \begin{equation}
    \vcenter{\infer[\constructor{(-)_\bot\text{-}isSet}]{\term{p} \equiv \term{q}}{\term{p} , \term{q} : \term{x}~\relation{\sqsubseteq_{\bot}}~\term{y}}}
  \end{equation}
\end{defn}

\subsection{Delay monad to Sequences}
\todo{introduction to subsection}
\begin{defn}
  We define
  \begin{equation}
    \typeformer{Seq}_{\type{R}} = \sum_{(\function{s} : \mathbb{N} \rightarrow \type{R} + \unit)} \function{isMon}~\function{s}
  \end{equation}
  where
  \begin{equation}
    \function{isMon}~\function{s} = \prod_{(n : \mathbb{N})} (\function{s}_n \equiv \function{s}_{n+1})  + ((\function{s}_n \equiv \mathtt{inr}~\unitelem) \times (\function{s}_{n+1} \not\equiv \mathtt{inr}~\unitelem))
  \end{equation}
  meaning a sequences is \(\mathtt{inr}~\unitelem\) until it reaches a point where it switches to \(\mathtt{inl}~\term{r}\) for some value \(\term{r}\). There are also the special cases of already terminated, meaning only \(\mathtt{inl}~\term{r}\) and never teminating meaning only \(\mathtt{inr}~\unitelem\).
\end{defn}
\noindent For each index in a sequence, the element at that index \(\function{s}_n\) is either not terminated \(\function{s}_n \equiv \mathtt{inr}~\unitelem\), which we denote as \(\function{s}_n \uparrow_{\type{R} + \unit}\), or it is terminated \(\function{s}_n \equiv \mathtt{inl}~\term{r}\) with some value \(r\), denoted by \(\function{s}_n \downarrow_{\type{R} + \unit} \term{r}\) or just \(\function{s}_n \downarrow_{\type{R} + \unit}\) to mean \(\function{s}_n \not\equiv \mathtt{inr}~\unitelem\). Thus we can write \(\function{isMon}\) as
\begin{equation}
  \function{isMon}~\function{s} = \prod_{(n : \mathbb{N})} (\function{s}_n \equiv \function{s}_{n+1})  + ((\function{s}_n \uparrow_{\type{R} + \unit}) \times (\function{s}_{n+1} \downarrow_{\type{R} + \unit}))
\end{equation}
We also introduce notation for the two special cases of sequences given above
\begin{equation}
  \function{now}_{Seq}~\term{r} = (\lambda\,\_, \mathtt{inl}~\term{r}) , (\lambda\,\_, \mathtt{inl}~\mathtt{refl})
\end{equation}
\begin{equation}
  \constant{never}_{Seq} = (\lambda\,\_, \mathtt{inr}~\unitelem) , (\lambda\,\_, \mathtt{inl}~\mathtt{refl})
\end{equation}
\todo[inline]{Some comment about decidable equivalence needed to show that \(\function{s}_{n+1} \not\equiv \mathtt{inr}~\unitelem\)}
\begin{defn}
  We can shift a sequence \((\function{s},\function{q})\) by inserting an element (and an equality) \((\term{z_s},\term{z_q})\) at \(n = 0\),
  \begin{equation}
    \function{shift}~(\function{s} , \function{q})~(\term{z_s} , \term{z_q}) = \begin{cases} \term{z_s} & n = 0 \\ \function{s}_m & n = m + 1 \end{cases} , \begin{cases} \term{z_q} & n = 0\\ \function{q}_m & n = m + 1 \end{cases},
  \end{equation}
\end{defn}
\begin{defn}
  We can unshift a sequence by removing the first element of the sequence
  \begin{equation}
    \begin{aligned}
      \function{unshift}&~(\function{s} , \function{q}) = \function{s} \circ \constructor{succ} , \function{q} \circ \constructor{succ}.
    \end{aligned}
  \end{equation}
\end{defn}
\begin{lem}
  The function
  \begin{equation}
    \function{shift\text{-}unshift}~(\function{s},\function{q}) = \function{shift}~(\function{unshift}~(\function{s},\function{q}))~(\function{s}_0,\function{q}_0)
  \end{equation}
  is equal to the identity function.
  \begin{proof}
    Unshifting a value followed by a shift, where we reintroduce the value we just remove, gives the sequence we started  with.
  \end{proof}
\end{lem}
\begin{lem}
  The function
  \begin{equation}
    \function{unshift\text{-}shift}~(\function{s},\function{q})=\function{unshift}~(\function{shift}~(\function{s},\function{q})~\_)
  \end{equation}
  is equal to the identity function.
  \begin{proof}
    If we shift followed by an unshift, we just introduce a value to instantly remove it, meaning the value does not matter.
  \end{proof}
\end{lem}
\noindent We now define an equivalence between \(\typeformer{delay}~\type{R}\) and \(\typeformer{Seq}_{\type{R}}\), where \(\constructor{later}\) are equivalent to shifts, and \(\constructor{now}~\term{r}\) is equivalent terminated sequence with value \(\term{r}\). We do this by defining equivalence functions, and the left and right identities.
\begin{lem}[\(\constant{inl{\not\equiv}inr}\)]
  For any two elements \(\term{x} = \mathtt{inl}~\term{a}\) and \(\term{y} = \mathtt{inr}~\term{b}\) then \(\term{x} \not\equiv \term{y}\).
  \begin{proof}
    The constructors are disjoint, so there is not a path between them \todo{better formulated proof}.
  \end{proof}
f\end{lem}
\begin{defn}
  We define a function from \(\typeformer{Delay}~\type{R}\) to \(\typeformer{Seq}_{\type{R}}\)
  \begin{equation}
    \begin{aligned}
      \function{Delay{\rightarrow}Seq}&~(\constructor{now}~\term{r}) = \function{now}_{Seq}~\term{r} \\
      \function{Delay{\rightarrow}Seq}&~(\constructor{later}~\term{x}) = \\[-3mm]
      & \strut\hspace{-8mm}\function{shift}~(\function{Delay{\rightarrow}Seq}~\term{x})~\left(\mathtt{inr}~\unitelem , \begin{cases} \mathtt{inr}~(\mathtt{refl} , \constant{inl{\not\equiv}inr}) & \term{x} = \constructor{now}~\_ \\ \mathtt{inl}~\mathtt{refl} & \term{x} = \constructor{later}~\_ \end{cases}\right)
    \end{aligned}
  \end{equation}
\end{defn}
\begin{defn}
  We define function from \(\typeformer{Seq}_{\type{R}}\) to \(\typeformer{Delay}~\type{R}\)
    \begin{equation}
      \begin{aligned}
        \function{Seq{\rightarrow}Delay}&~(\function{s} , \function{q}) = \begin{cases} \constructor{now}~\term{r} & \function{s}_0 = \mathtt{inl}~\term{r} \\ \constructor{later}~(\function{Seq{\rightarrow}Delay}~(\function{unshift}~(\function{s} , \function{q}))) & \function{s}_0 = \mathtt{inr}~\unitelem \end{cases}
      \end{aligned}
    \end{equation}
\end{defn}
\begin{thm}
  The type \(\typeformer{Seq}_{\type{R}}\) is equal to \(\typeformer{Delay}~\type{R}\)
  \begin{proof}  
    We define right and left identity, saying that for any sequence \((\function{s},\function{q})\), we get
    \begin{equation}
      \function{Delay{\rightarrow}Seq}~(\function{Seq{\rightarrow}Delay}~(\function{s},\function{q})) \equiv (\function{s},\function{q})
    \end{equation}
    defined by cases analysis on \(\function{s}_0\), if \(\function{s}_0 = \mathtt{inl}~\term{r}\) then we need to show
    \begin{equation}
      \function{now}_{Seq}~\term{r} \equiv (\function{s},\function{q})
    \end{equation}
    This is true, since \((\function{s},\function{q})\) is a monotone sequence and \(\mathtt{inl}~\term{r}\) is the top element of the order, then all elements of the sequence are \(\mathtt{inl}~\term{r}\). If \(\function{s}_0 = \mathtt{inr}~\unitelem\) then, we need to show
    \begin{equation}
      \function{shift}~(\function{Delay{\rightarrow}Seq}~(\function{Seq{\rightarrow}Delay}~(\function{unshift}~(\function{s},\function{q}))))~(\mathtt{inr}~\unitelem,\_) \equiv (\function{s},\function{q})
    \end{equation}
    by the induction hypothesis we get
    \begin{equation}
      \function{Delay{\rightarrow}Seq}~(\function{Seq{\rightarrow}Delay}~(\function{unshift}~(\function{s},\function{q}))) \equiv \function{unshift}~(\function{s},\function{q})
    \end{equation}
    since shift and unshift are inverse, we get the needed equality.
    \todo[inline]{Shift takes two arguemnts, either clarify that its shift' that inserts inr tt or ...}
    For the left identity, we need to show that for any delay monad \(\term{t}\) we get
    \begin{equation}
      \function{Seq{\rightarrow}Delay}~(\function{Delay{\rightarrow}Seq}~\term{t}) \equiv \term{t}
    \end{equation}
    defined by case analysis on \(\term{t}\), if \(\term{t} = \constructor{now}~a\) then the equality is \(\mathtt{refl}\). If \(\term{t} = \constructor{later}~\term{x}\) then we need to show
    \begin{equation}
      \constructor{later}~(\function{Seq{\rightarrow}Delay}~(\function{unshift}~(\function{shift}~(\function{Delay{\rightarrow}Seq}~x))) \equiv \constructor{later}~x
    \end{equation}
    By unshift and shift being inverse, and the induction hypothesis we get the wanted equality. Since we are able to define a left and right identity function, we get the wanted equality.
  \end{proof}
\end{thm}
\begin{cor}
  The types \(\typeformer{Delay}/\sim\) and \(\typeformer{Seq}/\sim\) are equal.
  \begin{proof}
    We show that if \(a~\relation{\sim_{delay}}~b\) then \(\function{Delay{\rightarrow}Seq}~a~\relation{\sim_{Seq}}~\function{Delay{\rightarrow}Seq}~b\), \todo{Show this}
    \\ \\
    and we show if \(x~\relation{\sim_{Seq}}~y\) then \(\function{Seq{\rightarrow}Delay}~x~\relation{\sim_{Seq}}~\function{Seq{\rightarrow}Delay}~y\), \todo{Show this}
  \end{proof}
\end{cor}

\subsection{Sequence to Partiality Monad}
In this section we will show that assuming the "Axiom of Countable Choice", we get an equivalence between sequences and the partiality monad.
\begin{defn}[Sequence Termination]
  The following relations says that a sequence \((\function{s},\function{q}) : \typeformer{Seq}_{\type{R}}\) terminates with a given value \(\term{r} : \type{R}\),
  \begin{equation}
    (\function{s},\function{q})~\relation{\downarrow_{\mathtt{Seq}}}~\term{r} = \sum_{(n : \mathbb{N})} \function{s}_n \downarrow_{\type{R} + \unit} \term{r}.
  \end{equation}
\end{defn}
\begin{defn}[Sequence Ordering]
  \begin{equation}
    (\function{s},\function{q})~\relation{\sqsubseteq_{\mathtt{Seq}}}~(\function{t},\function{p}) = \prod_{(\term{a} : \type{R})} \left(\|\function{s}~\relation{\downarrow_{\mathtt{Seq}}}~\term{a}\| \rightarrow \| \function{t}~\relation{\downarrow_{\mathtt{Seq}}}~\term{a}\|\right)
  \end{equation}
  where \(\| \cdot \|\) is propositional truncation.
\end{defn}
\begin{defn} There is a conversion from \(\type{R} + \unit\) to the partiality monad \(\type{R}_\bot\)
  \begin{equation}
    \begin{aligned}
      &\function{Maybe{\rightarrow}(-)_\bot}~(\mathtt{inl}~\term{r}) = \constructor{\eta}~\term{r} \\
      &\function{Maybe{\rightarrow}(-)_\bot}~(\mathtt{inr}~\unitelem) = \constructor{\bot}
    \end{aligned}
  \end{equation}
\end{defn}
\begin{defn}[Maybe Ordering]
  \label{eq:maybe-order}
  Given some \(x,y : \type{R} + \unit\), the ordering relation is defined as
  \begin{equation}
    \term{x}~\relation{\sqsubseteq_{\type{R} + \unit}}~\term{y} = (\term{x} \equiv \term{y}) + ((\term{x}  \downarrow_{\type{R} + \unit}) \times (\term{y} \uparrow_{\type{R} + \unit}))
  \end{equation}
\end{defn}
\noindent This ordering definition is basically \(\function{isMon}\) at a specific index, so we can again rewrite \(\function{isMon}\) as
\begin{equation}
  \function{isMon}~\function{s} = \prod_{(n : \mathbb{N})} \function{s}_n~\relation{\sqsubseteq_{\type{R} + \unit}}~\function{s}_{n+1}
\end{equation}
This rewriting confirms that if \(\function{isMon}~\function{s}\), then \(\function{s}\) is monotone, and therefore a sequence of partial values\todo{there exists non-monotone sequences, it just follows our definition of a sequence.}.
\begin{lem}
  The function \(\function{Maybe{\rightarrow}(-)_\bot}\) is monotone, that is, if \(\term{x}~\relation{\sqsubseteq_{\type{A} + \unit}}~\term{y}\), for some \(\term{x}\) and \(\term{y}\), then \((\function{Maybe{\rightarrow}(-)_\bot}~\term{x})~\relation{\sqsubseteq_{\bot}}~(\function{Maybe{\rightarrow}(-)_\bot}~\term{y})\).
  \begin{proof}
    We do the proof by case.
    \begin{equation}
      \begin{aligned}
        &\function{Maybe{\rightarrow}(-)_\bot\text{-}mono}~(\function{inl}~\pathterm{p}) = \\[-1mm]
        &\qquad\qquad \mathtt{subst}~(\lambda\,\term{a},~\function{Maybe{\rightarrow}(-)_\bot}~\term{x}~\relation{\sqsubseteq_{\bot}}~\function{Maybe{\rightarrow}(-)_\bot}~\term{a})~\pathterm{p}~(\constructor{\sqsubseteq_{refl}}~(\function{Maybe{\rightarrow}(-)_\bot}~\term{x})) \\[1mm]
        &\function{Maybe{\rightarrow}(-)_\bot\text{-}mono}~(\function{inr}~(\pathterm{p}, \_)) = \\[-1mm]
        &\qquad\qquad \mathtt{subst}~(\lambda\,\term{a},~\function{Maybe{\rightarrow}(-)_\bot}~\term{a}~\relation{\sqsubseteq_{\bot}}~\function{Maybe{\rightarrow}(-)_\bot}~\term{y})~\sym{\pathterm{p}}~(\constructor{\sqsubseteq_{never}}~(\function{Maybe{\rightarrow}(-)_\bot}~\term{y}))
      \end{aligned}
    \end{equation}
  \end{proof}
\end{lem}
\begin{defn}
  There is a function taking a sequence to an increasing sequence \todo{What is an increasing sequence ??, this is not defined anywhere!!}
  \begin{equation}
    \begin{aligned}
      &\function{Seq{\rightarrow}incSeq} \\
      &\function{Seq{\rightarrow}incSeq}~(\function{g} , \function{q}) = \function{Maybe{\rightarrow}(-)_\bot} \circ \function{g} , \function{Maybe{\rightarrow}(-)_\bot\text{-}mono} \circ \function{q}
    \end{aligned}
  \end{equation}
\end{defn}
\begin{defn}
  There is a function taking a sequence to the partiality monad
  \begin{equation}
    \begin{aligned}
      &\function{Seq{\rightarrow}(-)_\bot} : \typeformer{Seq}_{\type{A}} \rightarrow \type{A}_{\bot} \\
      &\function{Seq{\rightarrow}(-)_\bot}~(\function{g},\function{q}) = \bigsqcup \circ~\function{Seq{\rightarrow}incSeq}
    \end{aligned}
  \end{equation}
\end{defn}
\begin{lem}
  The function \(\function{Seq{\rightarrow}(-)_\bot}\) is monotone. 
  \begin{equation}
    \function{Seq{\rightarrow}(-)_\bot\text{-}mono} : \typeformer{isSet} \type{A} \rightarrow (\term{x}~\term{y} : \typeformer{Seq}_{\type{A}}) \rightarrow \term{x}~\relation{\sqsubseteq_{seq}}~\term{y} \rightarrow \function{Seq{\rightarrow}(-)_\bot}~\term{x}~\relation{\sqsubseteq_{\bot}}~\function{Seq{\rightarrow}(-)_\bot}~\term{y}
\end{equation}
  \begin{proof}
    Given two sequences, if one is smaller than the another, then the least upper bounds of each sequence respect the ordering. \todo{should this be formalized entirely, or should there just be a comment about monotonicity? Does not seem relevant? (There is alot of work here..)}
  \end{proof}
\end{lem}
\begin{defn}
  If two sequences \(\term{x},\term{y}\) are weakly bisimular, then \(\function{Seq{\rightarrow}(-)_\bot}~\term{x} \equiv \function{Seq{\rightarrow}(-)_\bot}~\term{y}\)
\begin{equation}
  \function{Seq{\rightarrow}(-)_\bot\text{-}{\approx}{\rightarrow}{\equiv}}~\term{A_{set}}~\term{x}~\term{y}~(\term{p} , \term{q}) = \constructor{\alpha_{\bot}}~(\function{Seq{\rightarrow}(-)_\bot\text{-}mono}~\term{A_{set}}~\term{x}~\term{y}~\term{p})~(\function{Seq{\rightarrow}(-)_\bot\text{-}mono}~\term{A_{set}}~\term{y}~\term{x}~\term{q})
\end{equation}
\end{defn}
\begin{defn}[Recursor for Quotient]
  \todo{is this a recursor, and for what? The quotient?}
  For all sequences \(\term{x}, \term{y} : \typeformer{Seq}_{\type{A}}\), functions \(\function{f} : \type{A} \rightarrow \type{B}\) and relations \(\function{g} : \term{x}~\relation{R}~\term{y} \rightarrow \function{f}~\term{x} \equiv \function{f}~\term{y}\), then if \(\type{B}\) is a set \(\term{B_{set}} : \typeformer{isSet}~\type{B}\), we get a function \(\function{rec} : \type{A} / \relation{R} \rightarrow \type{B}\), defined by case as
  \begin{equation}
    \label{eq:quotient-recursor}
    \begin{aligned}
      &\function{rec}~\function{[}\,\term{z}\,\function{]} = \function{f}~\term{z} \\
      &\function{rec}~(\function{eq/}~\_~\_~\term{r}~i) = \function{g}~\term{r}~i \\
      &\function{rec}~(\function{squash/}~\term{a}~\term{b}~\pathterm{p}~\pathterm{q}~i~j) = \term{B_{set}}~(\function{rec}~\term{a})~(\function{rec}~\term{b})~(\mathtt{ap}~\function{rec}~\pathterm{p})~(\mathtt{ap}~\function{rec}~\pathterm{q})~i~j
    \end{aligned}
  \end{equation}
\end{defn}
This recursor allows us to lift the function \(\function{Seq{\rightarrow}(-)_\bot}\) to the quotient
\begin{defn}
  \label{eq:QuotientedSeqToPartialityMonad}
  We can define a function \(\function{Seq/{\sim}{\rightarrow}(-)_\bot}\) from \(\typeformer{Seq}_{\type{A}}\) to \(\type{A}_{\bot}\), where \(\term{A_{set}} : \typeformer{isSet}~\type{A}\) as
  \begin{equation}
    \function{Seq/{\sim}{\rightarrow}(-)_\bot} = \function{rec}~\function{Seq{\rightarrow}(-)_\bot}~(\function{Seq{\rightarrow}(-)_\bot\text{-}{\approx}{\rightarrow}{\equiv}}~\term{A_{set}})~\constructor{(-)_\bot\text{-}isSet}
  \end{equation} 
\end{defn}
\begin{lem}
  \label{eq:injectivity-lemma}
  Given two sequences \(\term{s}\) and \(\term{t}\), if \(\function{Seq{\rightarrow}(-)_\bot}~\term{s} \equiv \function{Seq{\rightarrow}(-)_\bot}~\term{t}\), then \(\term{s}~\relation{\sim_{seq}}~\term{t}\).
  \begin{proof}
    We can reduce the burden of the proof, since
    \begin{equation}
      \term{s}~\relation{\sim_{seq}}~\term{t} = \left(\prod_{(\term{r} : \type{R})} \| \term{x}~\relation{\downarrow_{seq}}~\term{r} \| \rightarrow \| \term{y}~\relation{\downarrow_{seq}}~\term{r} \|\right) \times \left(\prod_{(\term{r} : \type{R})} \| \term{y}~\relation{\downarrow_{seq}}~\term{r} \| \rightarrow \| \term{x}~\relation{\downarrow_{seq}}~\term{r} \|\right)
    \end{equation}
    so we can just show one part and get the other by symmetry. We assume \(\| \term{x}~\relation{\downarrow_{seq}}~\term{r} \|\), to show \(\| \term{y}~\relation{\downarrow_{seq}}~\term{r} \|\). By the mapping property of propositional truncation, we reduce the proof to defining a function \(\term{x}~\relation{\downarrow_{seq}}~\term{r} \rightarrow \term{y}~\relation{\downarrow_{seq}}~\term{r}\). Since \(\term{x}~\relation{\downarrow_{seq}}~\term{r}\), then \(\constructor{\eta}~\term{r}~\relation{\sqsubseteq_{\bot}}~\function{Seq{\rightarrow}(-)_\bot}~\term{x}\), but we have assumed \(\function{Seq{\rightarrow}(-)_\bot}~\term{x} \equiv \function{Seq{\rightarrow}(-)_\bot}~\term{y}\), so we get \(\constructor{\eta}~\term{r}~\relation{\sqsubseteq_{\bot}}~\function{Seq{\rightarrow}(-)_\bot}~\term{y}\), and thereby \(\term{y}~\relation{\downarrow_{seq}}~\term{r}\).
  \end{proof}
\end{lem}
\begin{lem}
  The function \(\function{Seq/{\sim}{\rightarrow}(-)_\bot}\) is injective.
  \begin{proof}
    We use propositional elimination of quotients \todo{Should this be formalized?}
    \begin{equation}
      \begin{aligned}
        \function{elimProp} :~& (\typeformer{B} :\typeformer{Seq}_{\type{R}} / \relation{\sim_{seq}} \rightarrow \universe{\mathcal{U}}) \rightarrow ((\term{x} : \typeformer{Seq}_{\type{R}} / \relation{\sim_{seq}}) \rightarrow \typeformer{isProp}~(\typeformer{B}~\term{x})) \\
        \rightarrow~& (\function{f} : (\term{a} : \typeformer{Seq}_{\type{R}}) \rightarrow \typeformer{B}~\function{[}~a~\function{]}) \rightarrow (\term{x} : \typeformer{Seq}_{\type{R}} / \relation{\sim_{seq}}) \rightarrow \typeformer{B}~\term{x}
      \end{aligned}
    \end{equation}
    to show the injectivity, meaning for all \(\term{x}~\term{y} : \typeformer{Seq}_{\type{R}} / \relation{\sim_{seq}}\) we get \(\function{Seq/{\sim}{\rightarrow}(-)_\bot}~\term{x} \equiv \function{Seq/{\sim}{\rightarrow}(-)_\bot}~\term{y} \rightarrow \term{x} \equiv \term{y}\). We start by eliminating \(\term{x}\), followed by elimination of \(y\), this gives us the proof term \todo{Convert to text, instead of a proof term!?}
    \begin{equation}
      \begin{aligned}
        &\function{elimProp} \\
        &\quad (\lambda\,\term{a},~\function{Seq/{\sim}{\rightarrow}(-)_\bot}~\term{a} \equiv \function{Seq/{\sim}{\rightarrow}(-)_\bot}~\term{y} \rightarrow \term{a} \equiv \term{y}) \\
        &\quad (\lambda\,\term{a},~\function{isProp\Pi}~(\lambda\,\_, \constructor{squash/}~\term{a}~\term{y})) \\
        &\quad (\lambda\,\term{a},~\function{elimProp} \\
        &\qquad\qquad\quad (\lambda\,\term{b},~\function{Seq{\rightarrow}(-)_\bot}~\term{a} \equiv \function{Seq/{\sim}{\rightarrow}(-)_\bot}~\term{b} \rightarrow \function{[}~\term{a}~\function{]} \equiv \term{b}) \\
        &\qquad\qquad\quad (\lambda\,\term{b},~\function{isProp\Pi}~(\lambda\,\_, \constructor{squash/}~\function{[}~\term{a}~\function{]}~\term{b})) \\
        &\qquad\qquad\quad (\lambda\,\term{b},~(\constructor{eq/}~\term{a}~\term{b}) \circ (\function{Seq{\rightarrow}(-)_\bot\text{-}isInjective}~\term{a}~\term{b})))
      \end{aligned}
    \end{equation}
    where \(\function{Seq{\rightarrow}(-)_\bot\text{-}isInjective}\) is \eqref{eq:injectivity-lemma}, 
  \end{proof}
\end{lem}
\begin{lem}
  \label{eq:constant-sequence-equality}
  For all constant sequences \(\term{s}\), where all elements have the same value \(\term{v}\), we get \(\function{Seq{\rightarrow}(-)_\bot}~\term{s} \equiv \function{Maybe{\rightarrow}(-)_\bot}~\term{v}\).
  \begin{proof}
    The left side of the equality reduces to \(\function{Maybe{\rightarrow}(-)_\bot}\) applied on the least upper bound of the constant sequence, which is exactly the right hand side of the equality.
  \end{proof}
\end{lem}
\begin{lem}
  Assuming countable choice, the function \(\function{Seq{\rightarrow}(-)_\bot}\) is surjective
  \todo[inline]{describe countable choice (and why it is needed!)}
  \begin{proof}
    We do the proof by case \todo{describe what it means to do the surjective proof by case!} on \(\type{R}_{\bot}\), if it is \(\constructor{\eta}~r\) or \(\constructor{never}\), we convert them to the sequences \(\function{now}_{seq}~\term{r}\) and \(\constant{never}_{seq}\) respectively, then we are done by \eqref{eq:constant-sequence-equality}. For the least upper bound \(\constructor{\bigsqcup} (\function{s},\function{b})\), we translate to the (increasing) sequence, defined by \((\function{s},\function{b})\) \todo{more precise description!}. \todo{Complete the rest of the proof!}
  \end{proof}
\end{lem}
\begin{lem}
  Assuming countable choice, the function \(\function{Seq/{\sim}{\rightarrow}(-)_\bot}\) is surjective
  \begin{proof}
    \todo{Complete proof}
  \end{proof}
\end{lem}
\begin{thm}
  Assuming countable choice, we get an equivalence between sequences and the partiality monad.
  \begin{proof}
 The function \(\function{Seq/{\sim}{\rightarrow}(-)_\bot}\) is injective and surjective assuming countable choice, meaning we get an equivalence, since we are working in hSets.
\end{proof}
\end{thm}

\subsubsection{Building the weak bisimulation on the M-type as a M-type}
\todo[inline]{Is this possible? Yes! Should it be included?}

\subsubsection{Building the Partiality Monad as a limit (Dialgebra?)}
\todo[inline]{Is this possible?}

\section{Permutation Trees}
\todo[inline]{introduction to section}
\subsection{Binary Trees}
\todo[inline]{introduction to section}
\begin{defn}
  Binary trees \(\type{bT}\) are given as the \texttt{M}-type defined by the container
  \begin{equation}
    \left(\unit+\unit, \begin{cases} \empt & \mathtt{inl}~\unitelem\\ \unit + \unit & \mathtt{inr}~\unitelem \end{cases}\right)
  \end{equation}
\end{defn}
\noindent For which we get the constructors\\[-10mm]
\begin{center}
  \strut
  \hfill
  \begin{minipage}{0.20\linewidth}
    \begin{equation}
      \vcenter{\infer{\constructor{leaf} : \type{bT}}{\strut}}
    \end{equation}
  \end{minipage}
  \hfill
  \begin{minipage}{0.25\linewidth}
    \begin{equation}
      \vcenter{\infer{\constructor{node}~\function{f} : \type{bT}}{\function{f} : \unit + \unit \rightarrow \type{bT}}}
    \end{equation}
  \end{minipage}
  \hfill
  \strut
\end{center}
We want to define trees where the permutation does not matter, that is the following rule is true
\begin{equation}
  \vcenter{\infer[\constructor{mix}]{\constructor{node}~\function{f} \equiv \constructor{node}~(\function{f} \circ \function{g})}{\function{f} : \unit + \unit \rightarrow \type{bT} & \function{g} : \unit + \unit \rightarrow \unit + \unit & \mathtt{isIso}~\function{g}}}
\end{equation}
this type of tree we call  \(\type{bTp}\). We can either define this as a quotient \(\type{bT} / \relation{\sim}\), where
\begin{center}
  \strut
  \hfill
  \begin{minipage}{0.3\linewidth}
    \begin{equation}
      \vcenter{\infer[\constructor{\sim_{leaf}}]{\constructor{leaf}~\relation{\sim}~\constructor{leaf}}{\strut}}
    \end{equation}
  \end{minipage}
  \hfill
  \begin{minipage}{0.35\linewidth}
    \begin{equation}
      \vcenter{\infer[\constructor{\sim_{node}}]{\constructor{node}~\function{f}~\relation{\sim}~\constructor{node}~\function{g}}{\prod_{(\term{x} : \unit + \unit)} \function{f}~\term{x}~\relation{\sim}~\function{g}~\term{x}}}
    \end{equation}
  \end{minipage}
  \hfill
  \strut
\end{center}
and \(\constructor{\sim_mix}\) is \(\constructor{mix}\) with \(\sim\) replacing the equality. Another way to define the type is as a QIIT type where \(\constructor{mix}\) and \(\constructor{bTp\text{-}isSet}\) are additional constructors. We will show these two definitions are equal
\begin{thm}
  There is an equality \(\type{bT}/\relation{\sim} \equiv \type{bTp}\)
  \begin{proof}
    \begin{equation}
      \begin{aligned}
        &\function{bT{\rightarrow}bTp}~\constructor{leaf}_{\type{bT}} = \constructor{leaf}_{\type{bTp}} \\
        &\function{bT{\rightarrow}bTp}~(\constructor{node}_{\type{bT}}~\function{f}) = \constructor{node}_{\type{bTp}}~(\function{bT{\rightarrow}bTp} \circ \function{f})
      \end{aligned}
    \end{equation}
    This definition takes weakly bisimilar elements to equivalent elements
    \begin{equation}
      \begin{aligned}
        &\function{bT{\rightarrow}bTp\text{-}{\approx}{\rightarrow}{\equiv}}~(\constructor{\sim_{leaf}}) = \mathtt{refl} \\
        &\function{bT{\rightarrow}bTp\text{-}{\approx}{\rightarrow}{\equiv}}~(\constructor{\sim_{node}}~\function{p}) = \mathtt{ap}~\constructor{node}_{\type{bTp}}~(\function{bT{\rightarrow}bTp\text{-}{\approx}{\rightarrow}{\equiv}} \circ \function{p})
      \end{aligned}
    \end{equation}
    Then we lift the definition to the quotient
    \begin{equation}
      \begin{aligned}
        &\function{bT/{\sim}{\rightarrow}bTp} = \function{rec}~\function{bT{\rightarrow}bTp}~(\function{bT{\rightarrow}bTp\text{-}{\approx}{\rightarrow}{\equiv}}~\term{A_{set}})~\constructor{bTp\text{-}isSet}
      \end{aligned}
    \end{equation}
  \end{proof}
\end{thm}

\subsection{Silhouette Trees}
We start by defining an \(\type{R}\) valued \(\type{E}\) branching tree, as the \texttt{M}-type given by the following container
\begin{equation}
  \left(\type{R} + \unit , \begin{cases} \empt & \mathtt{inl}~\term{a} \\ \type{E} & \mathtt{inr}~\unitelem \end{cases}\right)
\end{equation}
We get the constructors
\begin{equation}
  \vcenter{\infer{\constructor{leaf}~\term{a} : \typeformer{tree}~\type{R}~\type{E}}{\term{a} : \type{R}}}
\end{equation}
\begin{equation}
  \vcenter{\infer{\constructor{node}~\function{k} : \typeformer{tree}~\type{R}~\type{E}}{\function{k} : \type{E} \rightarrow \typeformer{tree}~\type{R}~\type{E}}}
\end{equation}
Then we define the weak bisimilarity relation \(\relation{\sim_{tree}}\)
\begin{equation}
  \vcenter{\infer[\constructor{\sim_{leaf}}]{\constructor{leaf}~x~\relation{\sim_{tree}}~\constructor{leaf}~y}{\strut}}
\end{equation}
\begin{equation}
  \vcenter{\infer[\constructor{\sim_{node}}]{\constructor{node}~\function{k_1}~\relation{\sim_{tree}}~\constructor{node}~\function{k_2}}{\prod_{(v : \type{E})}~\function{k_1}~\term{v}~\relation{\sim_{tree}}~\function{k_2}~\term{v}}}
\end{equation}
This is enough to define, what we call, silhouette trees, which are trees quotiented by this notion of weak bisimilarity, namely \(\typeformer{tree} / \relation{\sim_{tree}}\). We can also construct this type directly as a QIIT, with type constructors \todo{add all needed constructors}
\begin{equation}
  \vcenter{\infer{\constructor{leaf}_{\typeformer{sTree}} : \typeformer{sTree}~\type{E}}{\strut}}
\end{equation}
\begin{equation}
  \vcenter{\infer{\constructor{node}_{\typeformer{sTree}}~\function{k} : \typeformer{sTree}~\type{E}}{\function{k} : \type{E} \rightarrow \typeformer{sTree}~\type{E}}}
\end{equation}
And the ordering relation \((\cdot~\relation{\sqsubseteq_{sTree}}~\cdot)\) of how "defined" the trees are by the constructors \todo{add all needed constructors}
\begin{equation}
  \vcenter{\infer{\constructor{\alpha}_{\typeformer{sTree}}~\term{x}~\term{y} : \typeformer{sTree}~\type{E}}{\term{x}~\relation{\sqsubseteq_{sTree}}~\term{y} & \term{y}~\relation{\sqsubseteq_{sTree}}~\term{x}}}
\end{equation}
\begin{equation}
  \vcenter{\infer{\mathop{\constructor{\bigsqcup}}_{(\function{e} : \mathbb{N} \rightarrow \type{E})}~(\function{s}~\function{e})}{\function{s} : (\mathbb{N} \rightarrow \type{E}) \rightarrow \typeformer{sTree}~\type{E}}}
\end{equation}

\subsubsection{From \(\typeformer{tree}\) to \(\typeformer{Seq}_{tree}\)}
We now want to show the equivalence between these two constructions, to do this we define an intermediate construction \(\typeformer{Seq}_{tree}\), where we get an ordering on the "definedness" of trees.
\begin{defn}
  We define monotone increasing sequences of trees as \todo{ordering is container ordering not maybe?}, all breanches are monotone increasing \todo{specify branches increasing?}.
  \begin{equation}
    \typeformer{Seq}_{tree} = \prod_{(\function{e} : \mathbb{N} \rightarrow \type{E})} \sum_{(\function{s} : \mathbb{N} \rightarrow \type{R} + \unit)} \prod_{(n : \mathbb{N})} \function{s}_n~\relation{\sqsubseteq_{\type{R}+\unit}}~\function{s}_{n+1}
  \end{equation}
  where \(\relation{\sqsubseteq_{\type{R}+\unit}}\) is similar to the relation defined at \eqref{eq:maybe-order} \todo{how does it differ? Constructors are equal type is different (trees instead of delay)}.
\end{defn}
\begin{defn}
  We define a function to shift a \(\typeformer{Seq}_{tree}\), it takes \(\function{f} : \type{E} \rightarrow \typeformer{Seq}_{tree}\) as an argument. We let \(\function{s'} = \function{f}~\function{e}_0~(\function{e} \circ \constructor{succ})\), then the definition is given as
  \begin{equation}
    \begin{aligned}
      &\function{shift\text{-}seq}~\function{f} = \\[-5mm]
      &\qquad \lambda\,\function{e}, \natcases{\mathtt{inr}~\unitelem}{\pi_1~\function{s'}} , \left(\lambda\,n, \begin{cases} \mathtt{inr}~(\mathtt{refl} , \constant{inl{\not\equiv}inr}) & n = 0 \land \pi_1~\function{s'}~0 = \mathtt{inl}~r \\ \mathtt{inl}~\mathtt{refl} & n = 0 \land \pi_1~\function{s'}~0 = \mathtt{inr}~\unitelem \\ \pi_2~\function{s'}~m & n = m + 1 \end{cases}\right)
    \end{aligned}
  \end{equation}
\end{defn}
\begin{defn}
  We define a function to unshift a \(Seq_{tree}\)
  \begin{equation}
    \begin{aligned}
      &\function{unshift\text{-}seq}~\function{s}~\term{v} = \lambda\,\function{e},~(\pi_1~(\function{s}~(\natcases{\term{v}}{\function{e}})) \circ \constructor{succ}) , (\pi_2~(\function{s}~(\natcases{v}{\function{e}})) \circ \constructor{succ})
    \end{aligned}
  \end{equation}
\end{defn}
\begin{lem}
  Shift and unshift er inverse to each other
  \begin{proof}
    The same reasoning as for \todo{shift - unshift}
  \end{proof}
\end{lem}
\begin{defn}
  We get a function from trees to monotone sequences
  \begin{equation}
    \begin{aligned}
      &\function{tree{\rightarrow}Seq}~(\constructor{leaf}~\term{r}) = \lambda\,\_, (\lambda\,\_, \mathtt{inl}~\term{r}) , (\lambda\,\_, \mathtt{inl}~\mathtt{refl}) \\
      &\function{tree{\rightarrow}Seq}~(\constructor{node}~\function{k}) = \function{shift}~(\function{tree{\rightarrow}Seq} \circ \function{k})
    \end{aligned}
  \end{equation}
\end{defn}
\begin{defn}
  We get a function from monotone sequences to trees
  \begin{equation}
    \begin{aligned}
      &\function{Seq{\rightarrow}tree}~\function{s} = \begin{cases} \constructor{leaf}~\term{r} & \prod_{(\function{e} : \mathbb{N} \rightarrow \type{E})} \pi_1~(\function{s}~\function{e})~0 = \mathtt{inl}~\term{r} \\ \constructor{node}~(\function{Seq{\rightarrow}tree} \circ \function{unshift}~\function{s}) & o.w.  \end{cases}
    \end{aligned}
  \end{equation}
\end{defn}
\begin{lem}
  \label{lem:seq-first-is-leaf}
  If the first element in the sequence is terminated / a leaf, then the rest of the elements will also be terminated.
  \begin{equation}
    \left(\prod_{e : \mathbb{N} \rightarrow \type{E}} \pi_1~(\function{s}~\function{e})~0 = \mathtt{inl}~\term{r}\right) \Leftrightarrow \left(\function{s} \equiv \lambda\,\_,(\lambda\,\_,\mathtt{inl}~\term{r}),(\lambda\,\_,\mathtt{inl}~\mathtt{refl})\right)
  \end{equation}
  \begin{proof}
    Since the sequence is monotone, and \(\mathtt{inl}~\term{r}\) is the top element of the order, if the first element is \(\mathtt{inl}~\term{r}\), then the sequence must be \(\lambda\,\_,(\lambda\,\_,\mathtt{inl}~\term{r}),(\lambda\,\_,\mathtt{inl}~\mathtt{refl})\). The other direction is trivial.
  \end{proof}
\end{lem}
\begin{thm}
  The types \(\typeformer{tree}\) and \(\typeformer{Seq}_{tree}\) are equal
  \begin{proof}
    We construct an isomorphism by the functions \(\function{tree{\rightarrow}Seq}\) and \(\function{Seq{\rightarrow}tree}\), with right inverse given by two cases, one where the first element in the sequence is \(\mathtt{inl}~\term{r}\), meaning representing a leaf with value \(r\), then we need to show that \(\function{s} \equiv \lambda\,\_,(\lambda\,\_,\mathtt{inl}~\term{r}),(\lambda\,\_,\mathtt{inl}~\mathtt{refl})\) which follows from \lemref{seq-first-is-leaf}. Otherwise we need to show that
    \begin{equation}
      \function{shift}~(\function{tree{\rightarrow}Seq} \circ \function{Seq{\rightarrow}tree} \circ \function{unshift}~\function{s}) \equiv \function{s}
    \end{equation}
    By induction we get
    \begin{equation}
      \function{tree{\rightarrow}Seq} \circ \function{Seq{\rightarrow}tree} \circ \function{unshift}~\function{s} \equiv \function{unshift}~\function{s}
    \end{equation}
    then by the right inverse of the equality between shift and unshift, we are done. For the left inverse we do case analysis, using induction and the left inverse of the equality between shift and unshift
    \begin{equation}
      \begin{aligned}
        &\function{tree\text{-}Seq}~(\constructor{leaf}~\term{r}) = \mathtt{refl} \\
        &\function{tree\text{-}Seq}~(\constructor{node}~\term{k}) = \function{unshift\text{-}shift}~(\function{tree{\rightarrow}Seq} \circ \function{k}) \ct \function{tree\text{-}Seq}~k
      \end{aligned}
    \end{equation}
  \end{proof}
\end{thm}
\noindent We start by defining some ordering relation on \(\typeformer{Seq}_{tree}\)
\begin{defn}[Sequence Termination]
  The following relations says that a branche \(\function{e} : \mathbb{N} \rightarrow \type{E}\) of a sequence \(\function{s} : \typeformer{Seq}_{tree}\) terminates at depth \(n\),
  \begin{equation}
    (\function{s}~\function{e})~\relation{\downarrow_{\typeformer{Seq}_{tree}}}~n = (\function{s}~\function{e}~n) \downarrow_{\type{R} + \unit}.
  \end{equation}
\end{defn}
We define weak bisimilarity relation for sequences \todo{Is this defined for partiality monad?}
\begin{defn}
  \begin{equation}
    \function{s}~\relation{\sim_{\typeformer{Seq}_{tree}}}~\function{t} = \prod_{(e : \type{E})} \prod_{(n : \mathbb{N})} \left(\|(\function{s}~\function{e})~\relation{\downarrow_{\mathtt{Seq}}}~\term{n}\| \longleftrightarrow \| (\function{t}~\function{e})~\relation{\downarrow_{\mathtt{Seq}}}~\term{n}\|\right)
  \end{equation}
  where \(\| \cdot \|\) is propositional truncation.
\end{defn}
\begin{cor}
  The types \(\typeformer{tree}/\relation{\sim_{tree}}\) and \(\typeformer{Seq}_{tree} / \relation{\sim_{\typeformer{Seq}_{tree}}}\) are equal.
  \begin{proof}
    We follow the same strategy as for \(\typeformer{Delay}/{\sim}\) and \(\typeformer{Seq}/{\sim}\)
  \end{proof}
\end{cor}

\subsubsection{\(\typeformer{Seq}_{tree}\) to \(\typeformer{sTree}\)}
\todo{Introduction to Seq to tree}
\begin{defn}
  We define a function converting a sequence on trees to a monotone sequence on \typeformer{sTree}'s
  \begin{equation}
    \function{Seq{\rightarrow}incSeq}
  \end{equation}
\end{defn}
\begin{defn}
  There is a function from \(\typeformer{Seq}_{tree}\) to \(\typeformer{sTree}\)
  \begin{equation}
    \begin{aligned}
      &\function{Seq{\rightarrow}sTree}~\function{s} = \mathop{\constructor{\bigsqcup}}_{(\function{e} : \mathbb{N} \rightarrow \type{E})}~(\function{Seq{\rightarrow}incSeq}~\function{s}~\function{e})
    \end{aligned}
  \end{equation}
\end{defn}
\begin{lem}
  Given \(a~\relation{\sqsubseteq_{\typeformer{Seq}_{tree}}}~b\), then \(\function{Seq{\rightarrow}sTree}~a~\relation{\sqsubseteq_{sTree}}~\function{Seq{\rightarrow}sTree}~b\).
  \begin{proof}
    \todo{Complete proof}.
  \end{proof}
\end{lem}
\noindent Since this definition is monotone, it can be lifted to the quotiented sequences.
\begin{lem}
  If two elements are weakly bisimilar, then they are equal as \function{sTrees} \todo{describe better}
  \begin{equation}
    \begin{aligned}
      &\function{Seq{\rightarrow}sTree\text{-}{\approx}{\rightarrow}{\equiv}}~\term{x}~\term{y}~(\function{p},\function{q}) = \\
      &\qquad\constructor{\alpha_{sTree}}~(\function{Seq{\rightarrow}sTree\text{-}mono}~\term{x}~\term{y}~\function{p})~(\function{Seq{\rightarrow}sTree\text{-}mono}~\term{y}~\term{x}~\function{q})
    \end{aligned}
  \end{equation}
\end{lem}
\begin{defn}
  Function from \(\typeformer{Seq/{\sim}}\) to \(\typeformer{sTree}\)
  \begin{equation}
    \function{Seq/{\sim}{\rightarrow}sTree} = \function{rec}~\function{Seq{\rightarrow}sTree}~\function{Seq{\rightarrow}sTree\text{-}{\approx}{\rightarrow}{\equiv}}~\constructor{sTree\text{-}isSet}
  \end{equation}
\end{defn}
\begin{lem}
  \(\function{Seq/{\sim}{\rightarrow}sTree}\) is injective
  \todo{TODO}
\end{lem}
\begin{lem}
  \(\function{Seq/{\sim}{\rightarrow}sTree}\) is surjective
  \todo{TODO}
\end{lem}
\begin{thm}
  \(\function{Seq/{\sim}{\rightarrow}sTree}\) is an equivalence.
  \todo{TODO}
\end{thm}

\subsection{QM-types}
We want to define sequences based on \texttt{M}-types, here are some examples
\begin{align}
  \typeformer{Seq}_{\typeformer{M}_{\containerpair{\type{A}}{\empt}}} &= \type{A} \\
  \typeformer{Seq}_{\typeformer{M}_{\containerpair{\type{A} + \unit}{[\empt,\unit]}}} &= \sum_{(\function{s} : \mathbb{N} \rightarrow \type{A} + \unit)} \prod_{(n : \mathbb{N})} \function{s}_n \sqsubseteq \function{s}_{n+1} \\
  \typeformer{Seq}_{\typeformer{M}_{\containerpair{\unit + \unit}{[\empt,\type{E}]}}} &= \prod_{(\function{e} : \mathbb{N} \rightarrow \type{E})}\sum_{(\function{s} : \mathbb{N} \rightarrow \unit + \unit)} \prod_{(n : \mathbb{N})} \function{s}_n \sqsubseteq \function{s}_{n+1} \\
  \typeformer{Seq}_{\typeformer{M}_{\containerpair{\type{A} + \unit}{[\empt,\type{E}]}}} &= \prod_{(\function{e} : \mathbb{N} \rightarrow \type{E})}\sum_{(\function{s} : \mathbb{N} \rightarrow \type{A} + \unit)} \prod_{(n : \mathbb{N})} \function{s}_n \sqsubseteq \function{s}_{n+1} \\
  \typeformer{Seq}_{\typeformer{M}_{\containerpair{\type{A} + \type{B}}{[\empt,\type{E}]}}} &= \prod_{(\function{e} : \mathbb{N} \rightarrow \type{E})}\sum_{(\function{s} : \mathbb{N} \rightarrow \type{A} + \type{B})} \prod_{(n : \mathbb{N})} \function{s}_n \sqsubseteq \function{s}_{n+1} \\
  \typeformer{Seq}_{\typeformer{M}_{\containerpair{\unit}{\unit + \unit}}} &=~\prod_{(\function{e} : \mathbb{N} \rightarrow \unit + \unit)}\sum_{(\function{s} : \mathbb{N} \rightarrow \type{1})}\prod_{(n : \mathbb{N})} \function{s}_n \sqsubseteq \function{s}_{n+1} \\
  \typeformer{Seq}_{\typeformer{M}_{\containerpair{\type{A}}{\unit + \unit}}} &=~\prod_{(\function{e} : \mathbb{N} \rightarrow \unit + \unit)}\sum_{(\function{s} : \mathbb{N} \rightarrow \type{A})}\prod_{(n : \mathbb{N})} \function{s}_n \sqsubseteq \function{s}_{n+1} \\
  \typeformer{Seq}_{\typeformer{M}_{\containerpair{\type{A}}{\typeformer{B}}}} &=~\prod_{(\function{e} : (n : \mathbb{N}) \rightarrow \typeformer{B}\,\function{s}_{n-1})}\sum_{(\function{s} : \mathbb{N} \rightarrow \type{A})}\prod_{(n : \mathbb{N})} \function{s}_n \sqsubseteq \function{s}_{n+1} \\
  \typeformer{Seq}_{\typeformer{M}_{\containerpair{\type{A}}{\typeformer{B}}}}^{(n)}~(\function{s}_{n-1} : \type{A}) &=~\sum_{(\function{s}_n : \type{A})} \prod_{(n : \mathbb{N})} \prod_{(\function{e} : (n : \mathbb{N}) \rightarrow \typeformer{B}\,\function{s}_{n-1})} \function{s}_n \sqsubseteq \function{s}_{n+1} \\
  \typeformer{Seq}_{\typeformer{M}_{\containerpair{\type{A} + \type{B}}{[\type{X} , \type{Y}]}}} &=~\prod_{\left(\function{e} : (n : \mathbb{N}) \rightarrow {\begin{cases}  {\scriptstyle  X~a} & {\scriptstyle  \function{s}_{n-1} = \mathtt{inl}~a} \\[-2mm] {\scriptstyle  Y~b} & {\scriptstyle  \function{s}_{n-1} = \mathtt{inr}~b} \end{cases}}\right)}\sum_{(\function{s} : \mathbb{N} \rightarrow \type{A})}\prod_{(n : \mathbb{N})} \function{s}_n \sqsubseteq \function{s}_{n+1}
\end{align}
\begin{equation}
  \typeformer{Seq}_{\typeformer{M}_{\containerpair{\type{A}}{\type{B}}}} = \sum_{(\term{s'} : \typeformer{M}_{\containerpair{\type{A}}{\typeformer{B}}})}  \sum_{(\function{s} : \mathbb{N} \rightarrow \sum_{(\term{a} : \type{A})} \typeformer{B}\,\term{a} \rightarrow \typeformer{M}_{\containerpair{\type{A}}{\typeformer{B}}})} \left(\prod_{\term{p} : \sum_{\term{a} : \type{A}} \typeformer{B}\,\term{a}} \term{s'} \sqsubseteq (\function{s}_0~\term{p})\right) \times \left( \prod_{(n : \mathbb{N})} \prod_{(p : \sum_{(a : A)} B\,a)} (s_n~p) \sqsubseteq (s_{n+1}~p) \right)
\end{equation}
\todo{which other QM types can be expressed as QIITs}
A QM-type is a quotiented M-type, we try to define this as a quotient on containers. We define container quotients as

\section{Cofree Coalgebra / Dialgebra}
\todo[inline]{Is this relevant?}

%%%%%%%%%%%%%%%%%%%%%%%%%%%%%%%%%%%%%%%%%%%%%%%%%%%%%%%%%%%%%%%%%%%%%%% 

\chapter{Properties of M-types?}
\section{Closure properties of \texttt{M}-types}
We want to show that \texttt{M}-types are closed under simple operations, we start by looking at the product.
\subsection{Product of \texttt{M}-types}
We start with containers and work up to \texttt{M}-types.
\begin{defn}
  The product of two containers is defined as \cite{ContainerRules}
  \begin{equation}
    \containerpair{\type{A}}{\typeformer{B}} \times \containerpair{\type{C}}{\typeformer{D}} \equiv \containerpair{\type{A} \times \type{C}}{\lambda\,(\term{a} , \term{c}), \typeformer{B}~\term{a} \times \typeformer{D}~\term{c}}.
  \end{equation}
\end{defn}
\noindent We can lift this rule, through the diagram in \figref{container-products}, used to define \texttt{M}-types.
\begin{figure}[h]
  \centering
  \begin{tikzcd}[row sep=huge,column sep=huge]
    \type{A} & \type{A} \times \type{C} \ar[l,"\pi_1"'] \ar[r,"\pi_2"] & \type{C} \\
    \functor{P} (\typeformer{M}_{\containerpair{\type{A}}{\typeformer{B}}}) \ar[u,"\pi_1"'] & \functor{P}(\typeformer{M}_{\containerpair{\type{A}}{\typeformer{B}} \times \containerpair{\type{C}}{\typeformer{D}}}) \ar[l] \ar[r] \ar[u,"\pi_1"'] & \functor{P} (\typeformer{M}_{\containerpair{\type{C}}{\typeformer{D}}}) \ar[u,"\pi_1"'] \\
    \typeformer{M}_{\containerpair{\type{A}}{\typeformer{B}}} \ar[u,"\function{out}"'] \ar[d,"\function{\pi}_{(!)}"] & \typeformer{M}_{\containerpair{\type{A}}{\typeformer{B}} \times \containerpair{\type{C}}{\typeformer{D}}} \ar[l] \ar[r] \ar[u,"\function{out}"'] \ar[d,"\function{\pi}_{(!)}"] & \typeformer{M}_{\containerpair{\type{C}}{\typeformer{D}}} \ar[u,"\function{out}"'] \ar[d,"\function{\pi}_{(!)}"] \\
    {\functor{P}_{\containerpair{\type{A}}{\typeformer{B}}}}^n~\unit \ar[d,"\function{\pi}_{(n)}"] & {\functor{P}_{\containerpair{\type{A}}{\typeformer{B}} \times \containerpair{\type{C}}{\typeformer{D}}}}^n~\unit \ar[l] \ar[r] \ar[d,"\function{\pi}_{(n)}"] & {\functor{P}_{\containerpair{\type{C}}{\typeformer{D}}}}^n~\unit \ar[d,"\function{\pi}_{(n)}"] \\
    \vdots \ar[d,"\function{\pi}_{(1)}"] & \vdots \ar[d,"\function{\pi}_{(1)}"] & \vdots \ar[d,"\function{\pi}_{(1)}"] \\
    \functor{P}_{\containerpair{\type{A}}{\typeformer{B}}}~\unit \ar[d,"\function{\pi}_{(0)}"] & \functor{P}_{\containerpair{\type{A}}{\typeformer{B}} \times \containerpair{\type{C}}{\typeformer{D}}}~\unit \ar[l] \ar[r] \ar[d,"\function{\pi}_{(0)}"] & \functor{P}_{\containerpair{\type{C}}{\typeformer{D}}}~\unit \ar[d,"\function{\pi}_{(0)}"] \\
    \unit & \unit \ar[l,"\texttt{refl}"] \ar[r,"\texttt{refl}"'] & \unit
  \end{tikzcd}
  \caption{Diagram for products of chains}
  \label{fig:container-products}
\end{figure}
\begin{thm}
  For any \(n : \mathbb{N}\) the following is true
  \begin{equation}
    {\functor{P}_{\containerpair{\type{A}}{\typeformer{B}}}}^n~\unit \times {\functor{P}_{\containerpair{\type{C}}{\typeformer{D}}}}^n~\unit \equiv {\functor{P}_{\containerpair{\type{A}}{\typeformer{B}} \times \containerpair{\type{C}}{\typeformer{D}}}}^n~\unit. \label{eq:polynomial-functor-product}
  \end{equation}
  
  \begin{proof}
    We do induction on \(n\), for \(n = 0\), we have \(\unit \times \unit \equiv \unit\). For \(n = m + 1\), we may assume
    \begin{equation}
      {\functor{P}_{\containerpair{\type{A}}{\typeformer{B}}}}^m~\unit \times {\functor{P}_{\containerpair{\type{C}}{\typeformer{D}}}}^m~\unit \equiv {\functor{P}_{\containerpair{\type{A}}{\typeformer{B}} \times \containerpair{\type{C}}{\typeformer{D}}}}^m~\unit,
    \end{equation}
    in the following
    \begin{align}
      &~ {\functor{P}_{\containerpair{\type{A}}{\typeformer{B}}}}^{m+1}~\unit \times {\functor{P}_{\containerpair{\type{C}}{\typeformer{D}}}}^{m+1}~\unit \\
      \equiv&~ \functor{P}_{\containerpair{\type{A}}{\typeformer{B}}} ({\functor{P}_{\containerpair{\type{A}}{\typeformer{B}}}}^m~\unit) \times \functor{P}_{\containerpair{\type{C}}{\typeformer{D}}} ({\functor{P}_{\containerpair{\type{C}}{\typeformer{D}}}}^m~\unit) \\
      \equiv&~ \sum_{\term{a} : \type{A}} \typeformer{B}~\term{a} \rightarrow {\functor{P}_{\containerpair{\type{A}}{\typeformer{B}}}}^m~\unit \times \sum_{\term{c} : \type{C}} \typeformer{D}~\term{c} \rightarrow {\functor{P}_{\containerpair{\type{C}}{\typeformer{D}}}}^m~\unit \\
      \equiv&~ \sum_{\term{a} , \term{c} : \type{A} \times \type{C}} (\typeformer{B}~\term{a} \rightarrow {\functor{P}_{\containerpair{\type{A}}{\typeformer{B}}}}^m~\unit) \times (\typeformer{D}~\term{c} \rightarrow {\functor{P}_{\containerpair{\type{C}}{\typeformer{D}}}}^m~\unit) \\
      \equiv&~ \sum_{\term{a} , \term{c} : \type{A} \times \type{C}} \typeformer{B}~\term{a} \times \typeformer{D}~\term{c} \rightarrow {\functor{P}_{\containerpair{\type{A}}{\typeformer{B}}}}^m~\unit \times {\functor{P}_{\containerpair{\type{C}}{\typeformer{D}}}}^m~\unit \\
      \equiv&~ \sum_{\term{a} , \term{c} : \type{A} \times \type{C}} \typeformer{B}~\term{a} \times \typeformer{D}~\term{c} \rightarrow {\functor{P}_{\containerpair{\type{A}}{\typeformer{B}} \times \containerpair{\type{C}}{\typeformer{D}}}}^m~\unit \\
      \equiv&~ \functor{P}_{\containerpair{\type{A}}{\typeformer{B}} \times \containerpair{\type{C}}{\typeformer{D}}} ({\functor{P}_{\containerpair{\type{A}}{\typeformer{B}} \times \containerpair{\type{C}}{\typeformer{D}}}}^m~\unit) \\
      \equiv&~ {\functor{P}_{\containerpair{\type{A}}{\typeformer{B}} \times \containerpair{\type{C}}{\typeformer{D}}}}^{m+1}~\unit
    \end{align}
    taking the limit of \eqref{eq:polynomial-functor-product} we get
  \begin{equation}
    \typeformer{M}_{\containerpair{\type{A}}{\typeformer{B}}} \times \typeformer{M}_{\containerpair{\type{C}}{\typeformer{D}}} \equiv \typeformer{M}_{\containerpair{\type{A}}{\typeformer{B}} \times \containerpair{\type{C}}{\typeformer{D}}}.
  \end{equation}
  \end{proof}
\end{thm}
\begin{exmp} For streams we get
\begin{equation}
  \typeformer{stream}~\type{A} \times \typeformer{stream}~\type{B} \equiv \typeformer{M}_{\containerpair{\type{A}}{\lambda~\_,\unit}} \times \typeformer{M}_{\containerpair{\type{B}}{\lambda~\_,\unit}} \equiv \typeformer{M}_{\containerpair{\type{A}}{\lambda~\_,\unit} \times \containerpair{\type{B}}{\lambda~\_,\unit}} \equiv \typeformer{stream}~(\type{A} \times \type{B}) \label{eq:zip-equality}
\end{equation}
as expected. Transporting along \eqref{eq:zip-equality} gives us a definition for \function{zip}.
\end{exmp}

\subsection{Co-product}
\todo[inline]{Coproducts?}

\subsection{...}
\todo[inline]{The rest of the closures defined in "Categories of Containers" \cite{ContainerRules}}

%%%%%%%%%%%%%%%%%%%%%%%%%%%%%%%%%%%%%%%%%%%%%%%%%%%%%%%%%%%%%%%%%%%%%%% 

\chapter{Examples of \texttt{M}-types}
\section{TODO: Place these subsections}
\todo[inline]{What makes a relation a bisimulation? Is bisim and equality equal.}
\subsection{Identity Bisimulation}
Lets start with a simple example of a bisimulation namely the one given by the identity relation for any \texttt{M}-type.
\begin{lem}
  The identity relation \((\cdot \equiv \cdot)\) is a bisimulation for any final coalgebra \(\coalg{\typeformer{M}_{\container{S}}}{\function{out}}\) defined over an \texttt{M}-type.
  
  \begin{proof}
  We first define the function
  \begin{equation}
    \begin{aligned}
      \function{\alpha}_{\equiv} &: \type{\overline{\equiv}} \rightarrow \functor{P}(\type{\overline{\equiv}}) \\
      \function{\alpha}_{\equiv} (x , y) &:= \pi_1~(\function{out}~x)~,~(\lambda\,b, (\pi_2~(\function{out}~x)~b \,,\, \texttt{refl}_{(\pi_2\,(\function{out}\,x)\,b)}))
    \end{aligned}
  \end{equation}
  and the two projections
  \begin{align}
    \pi_1^{\type{\overline{\equiv}}} &= (\pi_1 , \texttt{funExt}~\lambda\,(a,b,r), \texttt{refl}_{\function{out}\,a}) \\
    \pi_2^{\type{\overline{\equiv}}} &= (\pi_2 , \texttt{funExt}~\lambda\,(a,b,r), \texttt{cong}_{\function{out}}~(\sym{r})).
  \end{align}
  This defines the bisimulation, given by the diagram in \figref{id-bisim}.
  \begin{figure}[h]
    \centering
    \begin{tikzcd}
      \coalg{\typeformer{M}}{\function{out}} & \coalg{\type{\overline{\equiv}}}{\function{\alpha}_{\equiv}} \ar[l,Rightarrow,"{\pi_1}^{\type{\overline{\equiv}}}"'] \ar[r,Rightarrow,"{\pi_2}^{\type{\overline{\equiv}}}"] & \coalg{\typeformer{M}}{\function{out}}
    \end{tikzcd}
    \caption{Identity bisimulation}
    \label{fig:id-bisim}
  \end{figure}
\end{proof}
\end{lem}
\subsection{Bisimulation of Streams}
\todo[inline]{TODO}
\subsection{Bisimulation of Delay Monad}
We want to define a strong bisimulation relation \(\relation{\sim_{delay}}\) for the delay monad, 
\begin{defn}
  The relation \(\relation{\sim_{delay}}\) is defined by the following rules
  \begin{equation}
    \vcenter{\infer[\mathtt{now}{\sim}]{\function{now}~\term{r}~\relation{\sim_{delay}}~\function{now}~\term{r} : \universe{\mathcal{U}}}{\type{R} : \universe{U} & \term{r} : \type{R}}}
  \end{equation}
  \begin{equation}
    \vcenter{\infer[\mathtt{later}{\sim}]{\function{later}~\term{t}~\relation{\sim_{delay}}~\function{later}~\term{u} : \universe{\mathcal{U}}}{\type{R} : \universe{\mathcal{U}} & \term{t} : \typeformer{delay}~\type{R} & \term{u} : \typeformer{delay}~\type{R} & \term{t}~\relation{\sim_{delay}}~\term{u} : \universe{\mathcal{U}}}}
  \end{equation}
\end{defn}
\begin{thm}
  The relation \(\relation{\sim_{delay}}\) is a bisimulation for \(\typeformer{delay}~\type{R}\).
  \begin{proof}
    First we define the function
    \begin{equation}
      \begin{aligned}
        \function{\alpha}_{\relation{\sim_{delay}}} &: \type{\overline{\sim_{delay}}} \rightarrow \functor{P} (\type{\overline{\sim_{delay}}}) \\
        \function{\alpha}_{\relation{\sim_{delay}}}~(\term{a} , \term{b} , \mathtt{now}{\sim}~\term{r}) &:= (\mathtt{inr}~\term{r} , \lambda\,(\,)) \\
        \function{\alpha}_{\relation{\sim_{delay}}}~(\term{a} , \term{b} , \mathtt{later}{\sim}~\term{x}~\term{y}~\term{q}) &:= (\mathtt{inl}~\unitelem , \lambda\,\_, (\term{x} , \term{y} , \term{q}))
      \end{aligned}
    \end{equation}
    then we define the projections
      \begin{align}
        \pi_1^{\type{\overline{\sim_{delay}}}} &= \left( \pi_1 ~~ , ~~ \texttt{funExt}~\lambda\,(\term{a},\term{b},p), \begin{cases} (\mathtt{inr}~\term{r} , \lambda\,(\,)) & p = \mathtt{now}{\sim}~\term{r} \\ (\mathtt{inl}~\unitelem , \lambda\,\_, x) & p = \mathtt{later}{\sim}~\term{x}~\term{y}~\term{q}\end{cases} \right) \\
        \pi_2^{\type{\overline{\sim_{delay}}}} &= \left( \pi_2 ~~ , ~~ \texttt{funExt}~\lambda\,(\term{a},\term{b},p), \begin{cases} (\mathtt{inr}~\term{r} , \lambda\,(\,)) & p = \mathtt{now}{\sim}~\term{r} \\ (\mathtt{inl}~\unitelem , \lambda\,\_, y) & p = \mathtt{later}{\sim}~\term{x}~\term{y}~\term{q}\end{cases} \right) \\
      \end{align}
      This defines the bisimulation, given by the diagram in \figref{delay-strong-bisim}.
      \begin{figure}[h]
        \centering
        \begin{tikzcd}[column sep = large]
          \coalg{\typeformer{delay}~\type{R}}{\function{out}} & \coalg{\type{\overline{\sim_{\texttt{delay}}}}}{\function{\alpha}_{\relation{\sim_{delay}}}} \ar[l,Rightarrow,"\pi_1^{\type{\overline{\sim_{delay}}}}"'] \ar[r,Rightarrow,"\pi_2^{\type{\overline{\sim_{delay}}}}"] & \coalg{\typeformer{delay}~\type{R}}{\function{out}}
        \end{tikzcd}
        \caption{Strong bisimulation for delay monad}
        \label{fig:delay-strong-bisim}
      \end{figure}
  \end{proof}
\end{thm}

\subsection{Bisimulation of ITrees}
We define our bisimulation coalgebra from the strong bisimulation relation \(\mathcal{R}\), defined by the following rules.
\begin{equation}
  \vcenter{\infer[\mathtt{EqRet}]{\mathtt{Ret}~\term{a} \cong \mathtt{Ret}~\term{b}}{\term{a},\term{b} : \type{R} & \term{a} \equiv_{\type{R}} \term{b}}}
\end{equation}
\begin{equation}
  \vcenter{\infer[\mathtt{EqTau}]{\mathtt{Tau}~\term{t} \cong \mathtt{Tau}~\term{u}}{\term{t},\term{u} : \mathtt{itree}~\typeformer{E}~\type{R} & \term{t} \cong \term{u}}}
\end{equation}
\begin{equation}
  \vcenter{\infer[\mathtt{EqVis}]{\mathtt{Vis}~\term{e}~\term{k_1} \cong \mathtt{Tau}~\term{e}~\term{k_2}}{\type{A} : \universe{\mathcal{U}} & \term{e} : \typeformer{E}~\type{A} & \term{k_1},\term{k_2} : \type{A} \rightarrow \mathtt{itree}~\typeformer{E}~\type{R} & \term{t} \cong \term{u}}}
\end{equation}
Now we just need to define \(\term{\alpha}_{\relation{\mathcal{R}}}\) \todo[inline]{define the \(\alpha_{\relation{R}}\) function}. Now we have a bisimulation relation, which is equivalent to equality, using what we showed in the previous section.

\subsection{Zip Function}
We want the diagram in \figref{zip-computation-rules} to commute,
\begin{figure}[h]
  \centering
  \begin{tikzcd}[row sep=huge,column sep=huge]
    & A \times B & \\
    \typeformer{stream}~\type{A} \times \typeformer{stream}~\type{B} \ar[ur,"\destructor{hd} \times \destructor{hd}"'] \ar[rr,"\function{zip}"] \ar[d,"\destructor{tl} \times \destructor{tl}"] & & \typeformer{stream}~(\type{A} \times \type{B}) \ar[ul,"\destructor{hd}"] \ar[d,"\destructor{tl}"] \\
    \typeformer{stream}~\type{A} \times \typeformer{stream}~\type{B} \ar[rr,"\function{zip}"] & & \typeformer{stream}~(\type{A} \times \type{B})
  \end{tikzcd}
  \caption{TODO}
  \label{fig:zip-computation-rules}
\end{figure}
meaning we get the computation rules
\begin{equation}
  (\destructor{hd} \times \destructor{hd}) \equiv \destructor{hd} \circ \function{zip}
\end{equation}
\begin{equation}
  \function{zip} \circ (\destructor{tl} \times \destructor{tl}) \equiv \destructor{tl} \circ \function{zip}
\end{equation}
we can define the zip function as we did in the end of the last section. Another way to define the zip function is more directly, using the following lifting property of \texttt{M}-types
\begin{equation}
  \begin{aligned}
    &\function{lift_M}~\left(x : \prod_{n : \mathbb{N}} (\type{A} \rightarrow {\functor{P}_{\container{S}}}^n~\unit)\right)~\left(u : \prod_{n : \mathbb{N}} (\type{A} \rightarrow \pi_n (x_{n+1} a) \equiv x_n a)\right) (a : \type{A}) : \function{M}~\container{S} := \\
    &\qquad\qquad (\lambda\,n, x~n~a) , (\lambda\,n~i, p~n~a~i).
  \end{aligned} \label{eq:M-lift}
\end{equation}
To use this definition, we first define some helper functions
\begin{equation}
  \function{zip}_{X}~n~(x , y) = \begin{cases} \unit & \mathtt{if}~n = 0 \\ (\destructor{hd}~x , \destructor{hd}~y) , (\lambda\,\_, \function{zip}_{X}~m~(\destructor{tl}~x,\destructor{tl}~y)) , & \mathtt{if}~n = m + 1 \end{cases}
\end{equation}
\begin{equation}
  \function{zip}_{\pi}~n~(x , y) = \begin{cases} \function{refl} & \mathtt{if}~n = 0 \\ \lambda\,i,(\destructor{hd}~x , \destructor{hd}~y) , (\lambda\,\_, \function{zip}_{\pi}~m~(\destructor{tl}~x,\destructor{tl}~y)~i) , & \mathtt{if}~n = m + 1 \end{cases},
\end{equation}
we can then define
\begin{equation}
  \function{zip}_{lift}~(x , y) := \function{lift_M}~\function{zip}_X~\function{zip}~(x,y).
\end{equation}
\subsubsection{Equality of Zip Definitions}
We would expect that the two definitions for zip are equal
\begin{align}
  \texttt{transport}_{?}~a &\equiv \function{zip}_{lift}~a \\
                           &\equiv \function{lift_M}~\function{zip}_X~\function{zip}_\pi~(x,y) \\
                           &\equiv (\lambda\,n, \function{zip}_X~n~(x,y)) , (\lambda\,n~i, \function{zip}_\pi~n~(x,y)~i)
\end{align}
zero case \(X\)
\begin{equation}
  \function{zip}_X~0~(x,y) \equiv \unit
\end{equation}
Successor case \(X\)
\begin{align} 
  \function{zip}_X~(m+1)~(x,y) &\equiv (\destructor{hd}~x, \destructor{hd}~y),(\lambda\,\_,\function{zip}_X~m~(\destructor{tl}~x,\destructor{tl}~y)) \\
                               &\equiv (\destructor{hd}~x, \destructor{hd}~y),(\lambda\,\_, {?}~(\destructor{tl}~a)) \\
                               &\equiv (\destructor{hd}~(\texttt{transport}_{?} a)),(\lambda\,\_, \texttt{transport}_{?} (\destructor{tl}~a)) \\
                               &\equiv \mathtt{transport}_?~a \\
\end{align}
Zero case \(\pi\): \((\lambda\,i, \function{zip}_\pi~0~(x,y)~i \equiv refl)\).
\begin{align}
  &\equiv () , (\lambda\,i, \function{zip}_\pi~0~(x,y)~i) \\
  &\equiv \unit , \function{refl} \\
\end{align}
successor case
\begin{align}
  &\equiv (\function{zip}_X~(m+1)~(x,y)) , (\lambda\,i, \function{zip}_\pi~(m+1)~(x,y)~i) \\
  &\equiv ((\destructor{hd}~x, \destructor{hd}~y),(\lambda\,\_,\function{zip}_X~m~(\destructor{tl}~x,\destructor{tl}~y))) , (\lambda\,i, (\destructor{hd}~x, \destructor{hd}~y),(\lambda\,\_,\function{zip}_\pi~m~(\destructor{tl}~x,\destructor{tl}~y)~i))
\end{align}
\todo[inline]{Complete this proof}

\subsection{Examples of Fixed Points}
\subsubsection{Zeros}
Let us try to define the zero stream, we do this by lifting the functions
\begin{equation}
  \function{const_X}~(n : \mathbb{N})~(c : \mathbb{N})  := \begin{cases} \unit & n = 0 \\ (c , \lambda\,\_, \function{const_X}~m~c) & n = m + 1 \end{cases}
\end{equation}
\begin{equation}
  \function{const_\pi}~(n : \mathbb{N})~(c : \mathbb{N}) := \begin{cases} \mathtt{refl} & n = 0 \\ \lambda\,i,(c , \lambda\,\_, \function{const_\pi}~m~c~i) & n = m + 1 \end{cases}
\end{equation}
to get the definition of zero stream
\begin{equation}
  \mathtt{zeros} := \function{lift_M}~\function{const_X}~\function{const_\pi}~0.
\end{equation}
We want to show that we get the expected properties, such as
\begin{align}
  \destructor{hd}~\mathtt{zeros} &\equiv 0 \\
  \destructor{tl}~\mathtt{zeros} &\equiv \mathtt{zeros}
\end{align}

\subsubsection{Spin}
We want to define spin, as being the fixed point \(\mathtt{spin} = \mathtt{later}~\mathtt{spin}\), so that is again a final coalgebra, but of a \texttt{M}-type (which is a final coalgebra)

\begin{figure}[h]
  \centering
  \begin{tikzcd}[row sep=huge,column sep=huge]
    \function{spin} \ar[d,"\function{later}"] \\
    \function{later}~\function{spin}
  \end{tikzcd}
  \caption{TODO}
\end{figure}
\noindent Since it is final, it also must be unique, meaning that there is just one program that spins forever, without returning a value, meaning every other program must return a value. If we just


%%%%%%%%%%%%%%%%%%%%%%%%%%%%%%%%%%%%%%%%%%%%%%%%%%%%%%%%%%%%%%%%%%%%%%% 

\chapter{Conclusion}
\label{ch:conclusion}

\todo[inline]{conclude on the problem statement from the introduction}

%%%%%%%%%%%%%%%%%%%%%%%%%%%%%%%%%%%%%%%%%%%%%%%%%%%%%%%%%%%%%%%%%%%%%%% 

\cleardoublepage
\addcontentsline{toc}{chapter}{Bibliography}
\bibliographystyle{plain} 
\bibliography{Thesis}

%%%%%%%%%%%%%%%%%%%%%%%%%%%%%%%%%%%%%%%%%%%%%%%%%%%%%%%%%%%%%%%%%%%%%%% 

\cleardoublepage
\appendix

\chapter{Additions to the Cubical Agda Library}
% Should be in the appendix ...

%%%%%%%%%%%%%%%%%%%%%%%%%%%%%%%%%%%%%%%%%%%%%%%%%%%%%%%%%%%%%%%%%%%%%%% 

\chapter{The Technical Details}

\todo[inline]{\dots}

\end{document}

% read / cite : QUOTIENT INDUCTIVE-INDUCTIVE TYPES (https://arxiv.org/pdf/1612.02346.pdf)
% Compare lean paper to this body of work.
% read / cite : Quotiented containers
% [Abbott et al. 2005]
% if a type theory has W-types [Abbott et al.2005], then it has allinductive types
% read / cite : Constructing Quotient Inductive-Inductive Types (https://akaposi.github.io/finitaryqiit.pdf)
% tt in tt, http://www.cs.nott.ac.uk/~psztxa/publ/tt-in-tt.pdf
% http://cs.ioc.ee/ewscs/2017/studsess/veltri-slides.pdf

% Buzz words / key words:
% generalised algebraic theory
% non-well founded set
% inductive definition (also equalities)

% read up on "synthetic" homotopy theory

% On academic writing : https://oac.cdlib.org/findaid/ark:/13030/kt287035zn/entire_text/
% Mathematical Writing (CS 209) Videorecordings (V0174). Dept. of Special Collections and University Archives, Stanford University Libraries, Stanford, Calif.
% http://askarov.net/on-writing-and-presenting.html